\documentclass[11pt]{article}
\usepackage[T1]{fontenc}
\usepackage{lmodern}
\usepackage[margin=0.92in]{geometry}
\usepackage{amsmath,amssymb,amsthm,mathtools}
\usepackage{microtype}
\usepackage{booktabs,longtable,array,tabularx}
\usepackage[dvipsnames]{xcolor}
\usepackage[hyphens]{url}
\usepackage{hyperref}
\usepackage{enumitem}
\usepackage{listings}
\usepackage{fancyvrb}
\usepackage{graphicx}
\hypersetup{colorlinks=true,linkcolor=MidnightBlue,citecolor=MidnightBlue,
 urlcolor=MidnightBlue,pdftitle={Catalan Numbers in Intuitionistic Primitive Recursive Arithmetic},
 pdfauthor={Research prepared for Vladimir Reshetnikov},
 pdfsubject={Constructive Catalan identities, bijections, arithmetic, Hankel determinants, asymptotics, and Lean plan}}
\newcommand{\N}{\mathbb N}
\newcommand{\Z}{\mathbb Z}
\newcommand{\Q}{\mathbb Q}
\newcommand{\R}{\mathbb R}
\newcommand{\Cat}{\mathsf C}
\newcommand{\Bin}{\mathsf B}
\newcommand{\ipra}{\ensuremath{\mathrm{iPRA}}}
\newcommand{\monus}{\mathbin{\mathop{\dot{-}}}}
\newcommand{\card}[1]{\left|#1\right|}
\newcommand{\ind}[1]{\mathbf 1_{#1}}
\newcommand{\coeff}[2]{[x^{#1}]\,#2}
\newcommand{\vp}[1]{v_{#1}}
\newcommand{\floor}[1]{\left\lfloor #1\right\rfloor}
\newcommand{\ceil}[1]{\left\lceil #1\right\rceil}
\newcommand{\abs}[1]{\left|#1\right|}
\newcommand{\Dyck}{\mathsf{Dyck}}
\newcommand{\Ball}{\mathsf{Ball}}
\pdfstringdefDisableCommands{\def\ipra{iPRA}\def\Cat{C}}

\newtheorem{theorem}{Theorem}[section]
\newtheorem{lemma}[theorem]{Lemma}
\newtheorem{proposition}[theorem]{Proposition}
\newtheorem{corollary}[theorem]{Corollary}
\theoremstyle{definition}
\newtheorem{definition}[theorem]{Definition}
\newtheorem{construction}[theorem]{Construction}
\theoremstyle{remark}
\newtheorem{remark}[theorem]{Remark}
\newtheorem{status}[theorem]{Formal status}
\numberwithin{equation}{section}
\setlength{\emergencystretch}{3em}
\setlength{\parskip}{3.5pt}
\setcounter{tocdepth}{2}
\allowdisplaybreaks[2]
\lstset{basicstyle=\ttfamily\small,breaklines=true,columns=fullflexible,frame=single,
  xleftmargin=0.3em,xrightmargin=0.3em,showstringspaces=false,
  literate={≐}{{$\doteq$}}1}

\begin{document}
\hypersetup{pageanchor=false}
\begin{titlepage}
\centering
\vspace*{1.0cm}
{\Large Constructive arithmetic research\par}
\vspace{1.0cm}
{\huge\bfseries Catalan Numbers in Intuitionistic\\[7pt]
Primitive Recursive Arithmetic\par}
\vspace{0.7cm}
{\Large Finite bijections, exact identities, valuations,\\[4pt]
Hankel determinants, effective asymptotics, and a Lean blueprint\par}
\vspace{1.0cm}
{\large Prepared for Vladimir Reshetnikov\par}
\vspace{0.25cm}
{\large November 16, 2026\par}
\vfill
\begin{minipage}{0.90\textwidth}
\textbf{Principal result.}
The standard Catalan formulas and recurrences, all four additional finite sums
printed by MathWorld, the product, all six proof routes displayed by Wikipedia, the principal
interpretations listed by Wikipedia and MathWorld, the generating-function
and Bessel identities, the boundary sum,
the random-walk statement, prime and parity assertions, the three Hankel
determinants, the moment representation, rational and Fuss--Catalan
extensions, the $k$-fold convolution, and the displayed asymptotic terms all
admit constructive proofs whose arithmetical content is primitive recursive.
Infinite and analytic notation is replaced by formal coefficient identities or
by rational Cauchy data with explicit primitive-recursive moduli.

\medskip
\textbf{Checked base theorem.}
The inspected ProveIt snapshot already contains a Lean-kernel-checked
object-language derivation of the denominator-cleared Catalan identity.  The
present package starts from that theorem and gives a dependency-complete route
to the larger development.  The new Lean files in this delivery are an
uncompiled scaffold and executable specification; they contain no \texttt{sorry}.
\end{minipage}
\vfill
{\small ProveIt source snapshot:\quad
\texttt{00bbe871268c453903499fb6afe07519e87d9631}\par}
\end{titlepage}
\hypersetup{pageanchor=true}
\pagenumbering{roman}

\begin{abstract}
We develop a large part of the elementary and intermediate theory of Catalan
numbers inside intuitionistic primitive recursive arithmetic (\ipra).  The
basic numerical function is the division-free difference
\(\Cat_n=\binom{2n}{n}\monus\binom{2n}{n+1}\).  A reflection bijection proves
that it counts Dyck words, while the already checked ProveIt theorem proves
\((n+1)\Cat_n=\binom{2n}{n}\) in the embedded calculus
\(\mathsf{Prf}\,\mathsf{PRAAx}\).  From this base we derive all explicit
factorial, double-factorial, generalized-binomial, terminating
hypergeometric, product, recurrence, and finite-sum formulas displayed on the
requested pages.  Primitive-recursive bijections cover Dyck paths, balanced
parentheses, binary and plane trees, parenthesizations, triangulations,
noncrossing matchings and partitions, two-row standard tableaux, restricted
sequences, stack-sortable permutations, and all six length-three pattern
classes.

Formal power series are treated coefficientwise.  Analytic statements use
rational Cauchy codes.  An exact telescoping identity gives
\(\sum_{j=0}^N\Cat_j/4^j\) together with a tail at most
\(2/\sqrt{N+2}\), so the boundary sum and random-walk recurrence have very
small \ipra{} proofs.  Legendre and Kummer arguments give a digit algorithm for
\(v_p(\Cat_n)\), the parity criterion, the base-five last-digit criterion, and
the fact that only \(\Cat_2=2\) and \(\Cat_3=5\) are prime.  Ballot-number
Gram factorizations prove the two unit Hankel determinants; Cauchy--Binet and
an alternating null vector prove the shifted determinant \(n+1\), and the
first two determinant sequences characterize Catalan numbers uniquely.
A source audit also gives finite counterexamples to two incorrect Coxeter
identifications on the requested Wikipedia revision and proves the corrected
type-$A$ formulas.  Finally, finite Wallis inequalities and an
Euler--Maclaurin remainder give a primitive-recursive modulus for the leading equivalent and a quantitative
version of the first three terms printed by MathWorld.  A module-by-module
Lean plan targets ProveIt's object language and explicitly separates already
checked material from future formalization.
\end{abstract}

\tableofcontents
\newpage
\pagenumbering{arabic}

\section{What it means to prove a Catalan statement in \ipra}
\label{sec:foundations}

\subsection{The ambient theory and the inspected ProveIt base}

The intended theory is first-order intuitionistic logic with equality, a
function symbol for every primitive-recursive function together with its
defining equations, successor axioms, and induction for quantifier-free
formulas.  This is the usual presentation of intuitionistic primitive
recursive arithmetic.  No appeal to a completed infinite set, an unrestricted
minimum principle, excluded middle for arbitrary formulas, or countable choice
is needed below.

In ProveIt's embedded calculus, an internal theorem has type
\[
  ([]:\operatorname{List}(\mathsf{Fm}\;\mathsf{PTm}\;0))
  \ \vdash_{\mathsf{PRAAx}}\ \varphi.
\]
The relevant checked modules already provide factorials, exact division with
remainder, binomial coefficients, bounded sums and products, order,
pairing, primitive recursion, and determinant infrastructure.  In particular,
\texttt{Combinatorics/BinomialIdentities.lean} contains the theorem
\texttt{catalan}, whose object-language conclusion is the universal closure of
\[
 (n+1)\left(\binom{2n}{n}\monus\binom{2n}{n+1}\right)
   =\binom{2n}{n}.
\tag{1.1}\label{eq:checked-catalan}
\]
The source indexes the checked open variable as $n=m+1$; the case $n=0$ is a
one-line computation.  Thus integrality is not merely asserted in the
metatheory: the denominator-cleared identity is already a theorem of the
object theory.

\subsection{Finite objects are bounded natural-number codes}

Every combinatorial object in this report is represented by a natural number
below a primitive-recursive bound.  A word of length $r$ over an alphabet of
size $a$ is a number below $a^r$; its $i$th letter is extracted by quotient and
remainder.  A finite sequence is encoded by pairing, a finite tree by a prefix
word, a permutation by a bounded list with a primitive-recursive distinctness
test, and a finite set partition by a restricted-growth word.  All predicates
used below are quantifier-free after bounded searches are compiled to sums or
products of characteristic functions.

\begin{definition}[Bounded cardinality]
If $P(n,z)$ is a primitive-recursive Boolean predicate and $z<B(n)$ is a
primitive-recursive bound, put
\[
  \#P(n)=\sum_{z<B(n)}[P(n,z)],
\]
where $[P]\in\{0,1\}$ is the Boolean truth value.  This is itself a
primitive-recursive function.
\end{definition}

\begin{lemma}[Finite bijection principle]
\label{lem:finite-bijection}
Suppose $P(n,z)$ and $Q(n,w)$ have bounds $B(n)$ and $D(n)$.  If
primitive-recursive maps $F$ and $G$ satisfy, by quantifier-free bounded
proofs,
\[
 P(n,z)\Rightarrow Q(n,F(n,z)),\qquad
 Q(n,w)\Rightarrow P(n,G(n,w)),
\]
and $G(F(z))=z$, $F(G(w))=w$ on the indicated domains, then
$\#P(n)=\#Q(n)$ is provable in \ipra.
\end{lemma}

\begin{proof}
For each $w<D(n)$, the indicator of $Q(n,w)$ equals the bounded sum of
$[P(n,z)\wedge F(n,z)=w]$ over $z<B(n)$: existence follows from $G$, and
uniqueness follows by applying $G$ to two preimages.  Sum this equality over
$w$, interchange two bounded sums, and use the fact that exactly one $w$ is
$F(n,z)$.  Every operation is a bounded sum, equality test, and
quantifier-free induction.  No choice function is inferred: $F$ and $G$ are
part of the data.
\end{proof}

A recurrence proof has the same status.  If a class of size $n+1$ decomposes
uniquely into pairs of classes of sizes $i$ and $n-i$, the coding functions
for decomposition and reconstruction give
\[
  A_{n+1}=\sum_{i\le n}A_iA_{n-i}
\]
as an equality of bounded sums.  The word ``uniquely'' always means that the
inverse equations are explicitly checked.

\subsection{Division-free arithmetic}

Fractions are not primitive objects of \ipra.  Thus
\(\Cat_n=\binom{2n}{n}/(n+1)\) is represented by the pair of assertions
\[
  (n+1)\Cat_n=\binom{2n}{n},\qquad n+1>0,
\]
not by an unguarded division symbol.  Rational identities are cleared of all
positive denominators.  Signed identities are represented by pairs of natural
numbers or by ProveIt's integer-difference type.  A congruence
$a\equiv b\pmod m$ is the equation $\operatorname{rem}(a,m)=
\operatorname{rem}(b,m)$, with the positive-modulus hypothesis made explicit.
This convention is used throughout the number-theoretic and asymptotic
sections.

\subsection{Formal series mean finite coefficient identities}

A formal series $A(x)=\sum_{n\ge0}a_nx^n$ is not an infinite set in the object
language.  The assertion $A=B$ abbreviates
\[
  \forall N\;\bigl(a_N=b_N\bigr),
\]
and multiplication is the primitive-recursive Cauchy coefficient
\((AB)_N=\sum_{i\le N}a_i b_{N-i}\).  A square root with constant term $1$ is
specified recursively by its coefficients, or avoided by a polynomial
identity.  Consequently
\[
 C(x)=1+xC(x)^2
\tag{1.2}\label{eq:formal-catalan}
\]
is exactly Segner's finite recurrence, one coefficient at a time.  Formal
substitution is admitted only with zero constant term, so every requested
coefficient depends on finitely many input coefficients.

\begin{lemma}[Uniqueness of the Catalan formal solution]
There is at most one sequence $a_0,a_1,\ldots$ satisfying $a_0=1$ and
$a_{n+1}=\sum_{i\le n}a_i a_{n-i}$.
\end{lemma}

\begin{proof}
Quantifier-free induction on $n$.  At stage $n+1$ the right side uses only
indices at most $n$, which agree by the induction hypothesis.
\end{proof}

\subsection{Constructive real and complex notation}

A real number is used only through a primitive-recursive rational Cauchy code:
a pair $a^-(q)\le a^+(q)$ with width at most $1/q$, monotone refinement, and
an explicit modulus.  Addition, multiplication on bounded intervals, square
root by integer bisection, and logarithm and exponential by alternating or
Taylor sums have primitive-recursive realizers.  One convenient code for
$\pi$ is Machin's identity
\[
 \frac\pi4=4\arctan\frac15-\arctan\frac1{239},
\]
where the alternating arctangent series gives an immediate rational error
bound.  Complex numbers are pairs of such codes.

Accordingly, a limit statement $u_n\to L$ means a displayed
primitive-recursive $N(q)$ such that
\[
 n\ge N(q)\Longrightarrow \abs{u_n-L}\le \frac1q,
\]
a rational inequality after interval endpoints and denominators are cleared.
An integral means upper and lower rational Riemann sums with a displayed mesh
modulus.  The endpoint square root in Section~\ref{sec:series} is handled by a
change of variables before the sums are formed.

\subsection{Three levels of formal status}

To prevent an ambient Lean theorem from being confused with an internal
arithmetic theorem, the report uses the following labels.

\begin{enumerate}[label=\textbf{\Alph*.},leftmargin=2.4em]
\item \textbf{Checked object-language theorem}: a term of
      \(\mathsf{Prf}\,\mathsf{PRAAx}\) already accepted by Lean's kernel.
\item \textbf{iPRA-ready finite proof}: all functions, bounds, maps, and
      induction measures are explicit; translation to the existing bounded
      arithmetic library is routine but has not been entered as a proof term.
\item \textbf{Effective analytic proof}: the finite algebra is \ipra-ready
      and every limiting operation comes with a primitive-recursive rational
      modulus; the real-code library still has to be connected to
      \(\mathsf{Prf}\,\mathsf{PRAAx}\).
\end{enumerate}

Only Equation~\eqref{eq:checked-catalan} and the imported ProveIt lemmas are
claimed here at level A.  The new source files supplied with this package are
explicit scaffolds, not a claim of completed kernel verification.

\section{The arithmetic core and the explicit formulas}
\label{sec:arithmetic}

\subsection{Reflection, integrality, and the central formula}

Let a balanced word of size $n$ be a word of $n$ letters $U$ and $n$ letters
$D$.  Its height after $j$ letters is the number of $U$'s minus the number of
$D$'s in the first $j$ positions.  A Dyck word has all heights nonnegative.
There are $\binom{2n}{n}$ balanced words.

\begin{theorem}[Andre reflection]
\label{thm:reflection}
The number of bad balanced words, those having a negative height, is
$\binom{2n}{n+1}$.  Hence
\[
 \Cat_n:=\binom{2n}{n}\monus\binom{2n}{n+1}
\tag{2.1}\label{eq:cat-difference}
\]
is the number of Dyck words of semilength $n$.
\end{theorem}

\begin{proof}
For a bad word let $j$ be its least position of height $-1$.  The bounded
search for $j$ is primitive recursive.  Interchange $U$ and $D$ in the first
$j$ positions.  The resulting word has $n+1$ letters $U$ and $n-1$ letters
$D$.  Conversely, a word with these multiplicities first reaches height $1$;
reflect the prefix through that first visit.  The two bounded algorithms are
inverse.  The target words can be chosen by the $n+1$ positions of their $U$'s,
so their number is $\binom{2n}{n+1}$.  The binomial inequality needed for
monus follows at the same time, or directly from the explicit factor relation
below.
\end{proof}

Pascal arithmetic gives
\[
 (n+1)\binom{2n}{n+1}=n\binom{2n}{n}.
\tag{2.2}\label{eq:adjacent-binomial}
\]
Subtracting after multiplication by $n+1$ proves the checked identity
\[
 (n+1)\Cat_n=\binom{2n}{n}.
\tag{2.3}\label{eq:cat-central-cleared}
\]
Thus the usual fraction is a harmless abbreviation for a proven exact
quotient.

\begin{corollary}[Equivalent closed forms]
For $n\ge0$, with each quotient read in denominator-cleared form,
\begin{align}
 \Cat_n
 &=\frac1{n+1}\binom{2n}{n}
  =\frac{(2n)!}{(n+1)!n!},                                      \tag{2.4}\label{eq:factorial}\\
 &=\frac1{2n+1}\binom{2n+1}{n},                                 \tag{2.5}\label{eq:odd-binomial}\\
 (n\ge1)\qquad
 \Cat_n&=\frac1n\binom{2n}{n-1},                                \tag{2.6}\label{eq:n-binomial}\\
 (n+1)!\Cat_n&=2^n(2n-1)!!=\prod_{j=1}^{n}(4j-2).                \tag{2.7}\label{eq:product}
\end{align}
Here $(-1)!!=1$.
\end{corollary}

\begin{proof}
The factorial equation is the checked exactness theorem for binomial
coefficients followed by Equation~\eqref{eq:cat-central-cleared}.  The two
other binomial forms follow from
\[
 \binom{2n+1}{n}=\frac{2n+1}{n+1}\binom{2n}{n},\qquad
 \binom{2n}{n-1}=\frac n{n+1}\binom{2n}{n}.
\]
For the product, split $(2n)!$ into even and odd factors:
$(2n)!=2^n n!(2n-1)!!$, then multiply by $(n+1)!$.  Finally
$2(2j-1)=4j-2$ term by term.  All displayed fraction equations can be
replaced by cross multiplication, so the proof is quantifier-free algebra
after the imported factorial and binomial lemmas.
\end{proof}

The generalized-binomial formula printed by MathWorld is
\[
 \Cat_n=(-1)^n2^{2n+1}\binom{1/2}{n+1}.                          \tag{2.8}\label{eq:half-binomial}
\]
It is a rational identity rather than a new analytic theorem.  Expanding the
generalized binomial coefficient gives
\[
 \binom{1/2}{n+1}
  =\frac{(1/2)(-1/2)(-3/2)\cdots(1/2-n)}{(n+1)!}
  =\frac{(-1)^n(2n-1)!!}{2^{n+1}(n+1)!},
\]
and Equation~\eqref{eq:half-binomial} reduces to
Equation~\eqref{eq:product}.  Clearing $2^{n+1}(n+1)!$ makes this an integer
identity in \ipra.

The gamma expression
\[
 \Cat_n=\frac{4^n\Gamma(n+1/2)}{\sqrt\pi\,\Gamma(n+2)}             \tag{2.9}\label{eq:gamma}
\]
has the same finite content.  The half-integer recurrence
$\Gamma(x+1)=x\Gamma(x)$ and $\Gamma(1/2)=\sqrt\pi$ gives
\(\Gamma(n+1/2)=(2n)!\sqrt\pi/(4^n n!)\); substituting and cancelling
returns Equation~\eqref{eq:factorial}.  In an \ipra{} real-code library this is
a theorem about the indicated Cauchy names, but no general theory of the gamma
function is needed for integer $n$.

\subsection{Ratio and Segner recurrences}

\begin{theorem}[First-order recurrence]
\label{thm:ratio-recurrence}
For every $n\ge0$,
\[
 (n+2)\Cat_{n+1}=2(2n+1)\Cat_n.                                  \tag{2.10}\label{eq:ratio-cleared}
\]
Equivalently, since $n+2>0$,
\(\Cat_{n+1}/\Cat_n=2(2n+1)/(n+2)\).
\end{theorem}

\begin{proof}
Use Equation~\eqref{eq:factorial} twice and cancel the positive common factor
$(2n)!/(n!)^2$ in denominator-cleared form.  Alternatively, use the adjacent
central-binomial identity
$(n+1)\binom{2n+2}{n+1}=2(2n+1)\binom{2n}{n}$ and then
Equation~\eqref{eq:cat-central-cleared}.
\end{proof}

\begin{theorem}[Segner recurrence]
\label{thm:segner}
\[
 \Cat_0=1,
 \qquad
 \Cat_{n+1}=\sum_{i=0}^{n}\Cat_i\Cat_{n-i}.                      \tag{2.11}\label{eq:segner}
\]
\end{theorem}

\begin{proof}
Every nonempty Dyck word has a unique first-return decomposition
\[
 U\,u\,D\,v,
\]
where $u$ and $v$ are Dyck words.  If $u$ has semilength $i$, then $v$ has
semilength $n-i$.  The first return is found by bounded search; concatenation
is the primitive-recursive inverse.  Lemma~\ref{lem:finite-bijection} gives the
sum.  This proof is intuitionistic because it computes the decomposition; it
does not merely infer that some splitting point exists.
\end{proof}

Equations~\eqref{eq:ratio-cleared} and \eqref{eq:segner} each determine the
sequence from $\Cat_0=1$.  Their agreement can therefore also be proved by
induction after either recurrence has been established from the closed form.
The first values are
\[
 1,1,2,5,14,42,132,429,1430,4862,16796,58786.
\]
The accompanying exact Python checker evaluates these from all three
definitions.

\subsubsection{The three remaining finite proof variants on the Wikipedia page}

The first, second, and sixth proofs on the source page are respectively the
generating-function proof of Section~\ref{sec:series}, the reflection proof
above, and the cycle-lemma proof of Section~\ref{sec:generalizations}.  The
other three also have direct bounded implementations.

\paragraph{Chung--Feller exceedance classes.}
View a balanced path as a word in east and north steps from $(0,0)$ to
$(n,n)$, and let $e(P)$ be the number of north steps lying strictly above the
diagonal.  For $e(P)>0$, scan to the first north step entering the region
above the diagonal, then to the first east step $X$ that returns to the
diagonal.  If the step word factors as $F X L$, perform the cyclic cut
\[
                  F X L\longmapsto L X F.                       \tag{2.11a}
\]
All pieces are merely translated.  A height calculation at the two cut
points shows that precisely the distinguished north step ceases to be above
the diagonal, so $e$ decreases by one.  In the image, $X$ is the last east
step beginning on the diagonal; this property locates the inverse cut by a
bounded backward scan.  Iterating gives explicit bijections between every two
classes $e=r$ and $e=s$, for $0\le r,s\le n$.  The $n+1$ classes partition the
$\binom{2n}{n}$ balanced paths, and the class $e=0$ is the Dyck class.  Hence
$(n+1)\Cat_n=\binom{2n}{n}$.  This is the Chung--Feller proof in a form whose
forward map, inverse map, and iteration bound are primitive recursive.

\paragraph{The marked-triangulation ratio bijection.}
A triangulation of a convex $(n+2)$-gon has $(n+2)+(n-1)=2n+1$ edges, including
boundary edges and diagonals.  Thus triangulations with a marked oriented edge
number $(4n+2)\Cat_n$.  On the other hand, fix a base side of an
$(n+3)$-gon.  Its triangulations with one of the other $n+2$ boundary sides
marked number $(n+2)\Cat_{n+1}$.  There is a local inverse pair.  Open the
marked oriented edge of the first triangulated disk, insert one new boundary
vertex, and fill the opening by a triangle; mark the new non-base side.
Conversely, collapse the triangle incident with the marked non-base side and
remember the orientation induced on the collapsed edge.  On finite planar
codes, ``open'' and ``collapse'' are bounded list operations, and the two
operations are visibly inverse.  Therefore
\[
            (4n+2)\Cat_n=(n+2)\Cat_{n+1},                       \tag{2.11b}
\]
which is exactly Equation~\eqref{eq:ratio-cleared}.

\paragraph{Balanced strings and convolution cancellation.}
Write $B_n=\binom{2n}{n}$ for all balanced parenthesis strings of size $n$.
The first return of a nonempty balanced string is uniquely either
$UuDv$ or $D\bar uUv$, where $u$ is Dyck, $\bar u$ is obtained by swapping
$U,D$, and $v$ is balanced.  Consequently
\[
 B_{n+1}=2\sum_{j=0}^{n}\Cat_j B_{n-j}.                          \tag{2.11c}
\]
Subtracting Segner's recurrence gives
\[
 B_{n+1}-\Cat_{n+1}
   =\sum_{j=0}^{n}\Cat_j\bigl(2B_{n-j}-\Cat_{n-j}\bigr).         \tag{2.11d}
\]
There is a second count of the left side.  In a bad balanced string, stop at
the first down-step that reaches height $-1$.  The preceding word is a Dyck
word of some size $j$, and the remaining suffix has one more $U$ than $D$;
for $m=n-j$ it has $\binom{2m+1}{m}$ possible step words.  Hence
\[
 B_{n+1}-\Cat_{n+1}
   =\sum_{j=0}^{n}\Cat_j\binom{2(n-j)+1}{n-j}.                  \tag{2.11e}
\]
Since $\Cat_0=1$, coefficientwise cancellation of the two convolutions is a
plain induction: at stage $n$, every earlier summand is already equal, so the
last terms are equal.  Thus
\[
        2B_n-\Cat_n=\binom{2n+1}{n}.
\]
Solving for $\Cat_n$ and using
$\binom{2n+1}{n}=(2n+1)B_n/(n+1)$ yields
$\Cat_n=B_n/(n+1)$.  This completes constructive versions of all six proof
routes displayed on the Wikipedia page.

\subsection{The terminating hypergeometric and Narayana formulas}

For an integer $a$, let $(a)_k=a(a+1)\cdots(a+k-1)$, with $(a)_0=1$.
MathWorld's terminating hypergeometric expression is
\[
 {}_2F_1(1-n,-n;2;1)
 =\sum_{k=0}^{n}\frac{(1-n)_k(-n)_k}{(2)_k k!}.                   \tag{2.12}\label{eq:hypergeom}
\]
For $n\ge1$ the term $k=n$ vanishes.  For $0\le k<n$, cancellation of signed
factorials gives
\[
 \frac{(1-n)_k(-n)_k}{(2)_k k!}
 =\frac1{k+1}\binom{n-1}{k}\binom nk
 =\frac1n\binom n{k+1}\binom nk.                                \tag{2.13}\label{eq:narayana-term}
\]
Vandermonde's identity now yields
\[
 \sum_{k=0}^{n-1}\binom n{k+1}\binom nk=\binom{2n}{n-1},
\]
so Equations~\eqref{eq:n-binomial} and \eqref{eq:narayana-term} prove
\[
 {}_2F_1(1-n,-n;2;1)=\Cat_n.                                    \tag{2.14}
\]
The summands
\[
 N(n,j)=\frac1n\binom nj\binom n{j-1}\quad(1\le j\le n)
\]
are the Narayana numbers.  The proof above also proves their sum is $\Cat_n$.
All denominators divide the displayed integer numerators; the formal theorem
uses $nN(n,j)=\binom nj\binom n{j-1}$.

\subsection{The four additional finite sums on MathWorld}

\begin{theorem}[Two-colored Motzkin sum]
For $n\ge1$,
\[
 \Cat_n=\sum_{k=0}^{\floor{(n-1)/2}}
  \Cat_k\,2^{n-2k-1}\binom{n-1}{2k}.                             \tag{2.15}\label{eq:motzkin-sum}
\]
Terms with $2k>n-1$ are zero.
\end{theorem}

\begin{proof}
A two-colored Motzkin path of length $n-1$ consists of up, down, and two kinds
of level steps, never below height zero and ending at zero.  If it has $k$ up
steps, it has $k$ down steps.  Choose the $2k$ nonlevel positions, choose the
Dyck skeleton in $\Cat_k$ ways, and color the remaining positions in
$2^{n-1-2k}$ ways, giving the right side.

There is a primitive-recursive bijection from Dyck paths of semilength $n$:
remove the first $U$ and last $D$, group the remaining $2n-2$ letters into
adjacent pairs, and map
\[
 UU\mapsto U,\quad DD\mapsto D,\quad UD\mapsto L_0,\quad
 DU\mapsto L_1.
\]
Pair boundaries occur at even heights, so nonnegativity is preserved.  The
four substitutions are invertible, and restoring the outside letters gives
the inverse.
\end{proof}

\begin{theorem}[Weighted convolution]
For $n\ge1$,
\[
 \Cat_n=\frac1n\sum_{k=0}^{n-1}
 \Cat_{n-k-1}\binom{2k+1}{k+1}.                                 \tag{2.16}\label{eq:weighted-sum}
\]
\end{theorem}

\begin{proof}
Since $\binom{2k+1}{k+1}=(2k+1)\Cat_k$, the numerator is
\[
 \sum_{k=0}^{n-1}(2k+1)\Cat_k\Cat_{n-1-k}.
\]
Let $S=\sum k\Cat_k\Cat_{n-1-k}$.  Replacing $k$ by $n-1-k$ shows
$2S=(n-1)\Cat_n$ by Segner's recurrence.  Therefore the displayed numerator
is $2S+\Cat_n=n\Cat_n$.  The internal statement is this last equation, so no
division by a variable is required.
\end{proof}

\begin{theorem}[Alternating binomial transform]
For $n\ge0$,
\[
 \Cat_n=\sum_{k=0}^{n}(-1)^k2^{n-k}\binom nk
               \binom{k}{\floor{k/2}}.                           \tag{2.17}\label{eq:alternating-sum}
\]
\end{theorem}

\begin{proof}
Put $a_k=\binom{k}{\floor{k/2}}$ and
$A(t)=\sum_{k\ge0}a_kt^k$.  Separating even and odd indices gives, as a formal
coefficient identity,
\[
 A(t)=\frac{1+2t}{\sqrt{1-4t^2}}-tC(t^2),                        \tag{2.18}\label{eq:A-series}
\]
because
$\binom{2r+1}{r}=2\binom{2r}{r}-\Cat_r$.  Let the right side of
Equation~\eqref{eq:alternating-sum} be $b_n$ and
$B(z)=\sum b_nz^n$.  Interchanging two bounded coefficient sums gives
\[
 B(z)=\frac1{1-2z}A\left(\frac{-z}{1-2z}\right).                 \tag{2.19}\label{eq:Btransform}
\]
Use $\sqrt{1-4z}=1-2zC(z)$, itself just the polynomial relation
$C=1+zC^2$, and substitute Equation~\eqref{eq:A-series} into
Equation~\eqref{eq:Btransform}; after multiplication by $(1-2z)$ the result is
$(1-2z)C(z)$.  Thus $B=C$ coefficientwise.  Every coefficient extraction in
this calculation is a bounded sum, so the displayed formal-series proof is a
uniform \ipra{} proof.
\end{proof}

\begin{theorem}[Square sum]
For $n\ge0$,
\[
 \Cat_n=\sum_{k=0}^{\floor{n/2}}
 \left[\frac{n-2k+1}{n-k+1}\binom n{n-k}\right]^2.              \tag{2.20}\label{eq:square-sum}
\]
\end{theorem}

\begin{proof}
Adopt the convention $\binom nk=0$ outside $0\le k\le n$, and put
$d_k=\binom nk-\binom n{k-1}$.  For $k\le n/2$,
\[
 d_k=\frac{n-2k+1}{n-k+1}\binom nk.
\]
Symmetry gives $d_{n+1-k}=-d_k$, with a zero middle term when necessary.
Hence the sum in the statement is half of $\sum_{k=0}^{n+1}d_k^2$.  Expanding
and applying Vandermonde twice,
\begin{align*}
 \sum_{k=0}^{n+1}d_k^2
 &=2\sum_k\binom nk^2-2\sum_k\binom nk\binom n{k-1}\\
 &=2\binom{2n}{n}-2\binom{2n}{n-1}=2\Cat_n.
\end{align*}
This proof uses only integer identities.  Combinatorially, $d_k$ is the number
of standard tableaux of shape $(n-k,k)$, and the same sum follows from the
Robinson--Schensted correspondence for $123$-avoiding permutations; the
algebraic proof avoids making that larger bijection a dependency.
\end{proof}

\subsection{Riordan inverse transform}

Let $R_n$ be the number of noncrossing partitions of $[n]$ having no singleton
blocks, with $R_0=1$.  Every noncrossing partition is obtained uniquely by
choosing the elements that are not singleton blocks and placing a
singleton-free noncrossing partition on the chosen cyclically ordered subset.
Therefore
\[
 \Cat_n=\sum_{k=0}^{n}\binom nk R_k.                              \tag{2.21}\label{eq:riordan-transform}
\]
Binomial inversion gives
\[
 R_n=\sum_{k=0}^{n}(-1)^{n-k}\binom nk\Cat_k.                    \tag{2.22}
\]
Thus the two sequences are inverse under the binomial-transform operation.
The formulas themselves remove a convention trap.  MathWorld calls
\[
 b_n=\sum_{k=0}^n(-1)^{n-k}\binom nk a_k
 \quad\text{the binomial transform, and}\quad
 a_n=\sum_{k=0}^n\binom nk b_k
 \quad\text{its inverse}.
\]
With $a_n=\Cat_n$ and $b_n=R_n$, equations~\eqref{eq:riordan-transform}
and~(2.22) therefore say exactly that the Catalan numbers are the inverse
binomial transform of the Riordan numbers.  Under the alternative widespread
convention that calls the positive-sum operation the ``binomial transform,''
one instead says that Catalan is the binomial transform of Riordan.  The
deletion and insertion of singleton blocks are bounded primitive-recursive
inverse maps.

\section{Primitive-recursive Catalan bijections}
\label{sec:combinatorics}

This section proves the principal counting interpretations listed on the two
source pages.  Each map is an algorithm on bounded words or lists, and each
inverse is described.  Thus the arguments are not merely classical counting
proofs; they directly instantiate Lemma~\ref{lem:finite-bijection}.

\subsection{Words, parentheses, lattice paths, and mountains}

The Dyck-word interpretation is Theorem~\ref{thm:reflection}.  Replacing $U$
by an opening parenthesis and $D$ by a closing parenthesis gives correctly
matched parenthesis strings: the height is the number of currently open
parentheses.  Replacing $U,D$ by right and up steps gives monotone paths from
$(0,0)$ to $(n,n)$ that remain on the chosen side of the diagonal.  Replacing
them by slopes $+1,-1$ gives mountains that do not fall below the starting
level.  These are literal letter substitutions, so their inverses are
immediate.

The column-height representation on Wikipedia is also direct.  For the
$j$th right step, record the number $h_j$ of up steps preceding it.  A path
below the diagonal gives
\[
 0\le h_1\le h_2\le\cdots\le h_n,
 \qquad h_j\le j-1.
\tag{3.1}\label{eq:height-sequence}
\]
Conversely, insert the required number of up steps before each right step and
then the remaining up steps.  This gives a primitive-recursive bijection with
nondecreasing sequences satisfying Equation~\eqref{eq:height-sequence}.

\subsection{Full binary trees and complete parenthesizations}

A full ordered binary tree is either a leaf or an internal node with an ordered
left and right subtree.  Define its Dyck code recursively by
\[
 \operatorname{code}(\text{leaf})=\varepsilon,
 \qquad
 \operatorname{code}(\operatorname{node}(L,R))
   =U\operatorname{code}(L)D\operatorname{code}(R).              \tag{3.2}\label{eq:tree-code}
\]
If the tree has $n$ internal nodes, the word has $n$ $U$'s and $n$ $D$'s.
The first-return decomposition in Theorem~\ref{thm:segner} recovers $L$ and
$R$, so Equation~\eqref{eq:tree-code} is a bijection.  Hence $\Cat_n$ counts full
ordered binary trees with $n$ internal vertices and $n+1$ leaves.

A complete parenthesization of $n+1$ ordered factors has a unique outermost
binary operation; its two operands are parenthesized initial and final blocks.
Make this operation the root and recurse.  Reading a tree back as an operation
gives the inverse.  The same proof covers binary parse trees, bracketings of
ordered exponentials, and any fixed binary operator.  Adding a plant edge and
suppressing leaves converts a full binary tree into the planted planar
trivalent tree described by MathWorld.

\subsection{Rooted plane trees}

For a rooted ordered tree with $n$ edges, perform its contour traversal.  Write
$U$ when an edge is traversed away from the root for the first time and $D$
when it is traversed back.  The current depth is the Dyck height, so the output
is a Dyck word of semilength $n$.  Conversely, scan a Dyck word with a stack:
on $U$, create the next child of the current vertex and push it; on $D$, pop.
The prefix condition prevents underflow and the final height empties the
stack.  These two routines are inverse, proving that $\Cat_n$ counts rooted
plane trees with $n$ edges, equivalently $n+1$ vertices.

The familiar left-child/right-sibling representation gives the same
bijection through binary trees: the left pointer is the first child and the
right pointer is the next sibling.  This alternative is useful in Lean because
both trees can share one bounded-array representation.

\subsection{Triangulations of a convex polygon}

Fix a distinguished boundary edge of a convex $(n+2)$-gon.  In any
triangulation, exactly one triangle is incident with that edge.  If its third
vertex leaves $i+2$ vertices on one side and $n-i+2$ on the other, the remaining
diagonals are independent triangulations of those two polygons.  Thus the
numbers $T_n$ satisfy
\[
 T_0=1,\qquad T_{n+1}=\sum_{i=0}^{n}T_iT_{n-i}.                  \tag{3.3}
\]
More strongly, the incident triangle is a computable root, and recursively
placing the two subtriangulations to its left and right is a bijection with
full ordered binary trees.  Therefore $T_n=\Cat_n$.  The same recursion proves
Euler's polygon-division problem and Segner's formula
$E_m=\Cat_{m-2}$ for $m\ge2$.

\subsection{Noncrossing handshakes and pair partitions}

Place $2n$ labeled points around a circle.  In a noncrossing perfect matching,
point $1$ is paired with an even point $2i+2$.  The $2i$ points inside that
chord and the remaining $2(n-1-i)$ points outside it are matched
independently.  This is an explicit first-chord decomposition with inverse,
so the count satisfies Segner's recurrence and is $\Cat_n$.

Equivalently, scan the points in cyclic order and write $U$ at the first
endpoint and $D$ at the second endpoint of each chord, where chords are
processed by nesting.  Noncrossing is exactly the stack condition.  This also
proves the Wikipedia statement about noncrossing partitions of $[2n]$ whose
blocks all have size two.

\subsection{Noncrossing set partitions and the Bell bound}

Encode a partition of $[n]$ by its restricted-growth word.  It is noncrossing
when there do not exist $a<b<c<d$ with $a,c$ in one block and $b,d$ in another.
This is a bounded decidable predicate.

Let $P(x)$ be the formal counting series for noncrossing partitions.  In a
nonempty partition, list the elements of the block containing $1$ as
$1=b_0<b_1<\cdots<b_r$.  The $r+1$ cyclic intervals after these elements carry
independent noncrossing partitions.  Marking the $r+1$ elements of the root
block gives, coefficientwise,
\[
 P=1+\sum_{r\ge0}x^{r+1}P^{r+1}
  =1+\frac{xP}{1-xP}.                                           \tag{3.4}\label{eq:NCseries}
\]
Multiplication by $1-xP$ reduces this to $P=1+xP^2$.  By formal uniqueness,
$P=C$, so there are $\Cat_n$ noncrossing partitions of $[n]$.  The decomposition
and reconstruction are bounded list algorithms; the geometric terminology is
not essential to the proof.

The inclusion of noncrossing partitions into all set partitions is injective,
so
\[
 \Cat_n\le \operatorname{Bell}(n).                               \tag{3.5}\label{eq:bell-bound}
\]
This is an \ipra{} inequality between bounded counts, not an appeal to cardinal
comparison of abstract sets.

\subsection{Two-row standard Young tableaux}

Given a standard Young tableau of rectangular shape $(n,n)$, scan the numbers
$1,\ldots,2n$ and write $U$ if the number is in the first row and $D$ if it is
in the second.  Each row contains $n$ entries.  The $j$th first-row entry is
smaller than the $j$th second-row entry exactly when every prefix has at least
as many $U$'s as $D$'s.  Thus the word is Dyck.

Conversely, place the positions of the $U$'s increasingly in the first row and
the positions of the $D$'s increasingly in the second.  The Dyck prefix
condition is precisely the column inequality.  This proves that the number of
$2$-by-$n$ standard tableaux is $\Cat_n$.  Evaluating the hook-length product
also gives
\[
 \frac{(2n)!}{n!(n+1)!}=\Cat_n,
\]
so the hook-length specialization agrees with the bijection.

\subsection{The restricted sequences on Wikipedia}

Let $a_i$ be the height immediately after the $i$th $U$ in a Dyck word.  Then
\[
 a_1=1,
 \qquad 1\le a_{i+1}\le a_i+1.                                  \tag{3.6}\label{eq:restricted-seq}
\]
Thus the sequence may rise by at most one and may fall by any amount while
remaining positive.  Conversely, from such a sequence insert
$a_i-a_{i+1}+1$ letters $D$ between the $i$th and $(i+1)$st $U$, and insert
$a_n$ letters $D$ at the end.  These numbers are nonnegative and telescope to
$n$ down steps.  The construction stays nonnegative and inverts height
recording.  Hence these length-$n$ sequences are counted by $\Cat_n$.

\subsection{Pattern-avoiding and stack-sortable permutations}

A permutation avoids $231$ precisely when, after writing it as $u\,n\,v$ at
its largest entry, every entry of $u$ is smaller than every entry of $v$, and
both standardized subwords avoid $231$.  If some $x\in u$ exceeds some
$y\in v$, then $x,n,y$ is a $231$ pattern; the converse is immediate by
locating the largest member of any alleged pattern.  Therefore the left word
uses the smallest $i$ values for some $i$, the right word uses the remaining
values, and
\[
 A_{n+1}=\sum_{i=0}^{n}A_iA_{n-i},\qquad A_0=1.                  \tag{3.7}
\]
Standardization and unstandardization are primitive recursive, so
$A_n=\Cat_n$ by an actual decomposition bijection.

Reverse, complement, and inverse are explicit involutions on permutations.
They show that $132,213,231,312$ have the same avoidance count.  For the
remaining orbit, the Robinson--Schensted row-insertion algorithm is primitive
recursive and its reverse row-deletion algorithm is an explicit inverse.  The
special Schensted lemma needed here also has a bounded inductive proof: the
row-bumping lemma says that insertion paths for successively smaller letters
move strictly left and strictly downward.  Hence a decreasing subsequence of
length $r$ forces $r$ distinct rows.  Conversely, tracing the bumping
predecessors of the leftmost entry in each successive row, and then reversing
the recorded insertion times, constructs a decreasing subsequence with one
entry from every occupied row.  Thus a permutation avoids $321$ exactly when
its insertion tableau has at most two rows.  Therefore its number is
\[
 \sum_{k=0}^{\floor{n/2}} f_{(n-k,k)}^2,
\]
where the two-row hook calculation gives
$f_{(n-k,k)}=\binom nk-\binom n{k-1}$.  Equation~\eqref{eq:square-sum} identifies
this sum with $\Cat_n$.  Reversal covers $123$.  Consequently every pattern of
length three has exactly $\Cat_n$ avoiding permutations.

The recursive stack-sorting operator satisfies
$S(u\,n\,v)=S(u)S(v)n$.  Its output is increasing if and only if $S(u)$ and
$S(v)$ are increasing and every value in $u$ is smaller than every value in
$v$.  By the preceding characterization this is equivalent, inductively, to
$231$-avoidance.  Thus the stack-sortable permutations are also counted by
$\Cat_n$.

\subsection{Staircase tilings and related tree models}

A tiling of the staircase Ferrers shape $(n,n-1,\ldots,1)$ by $n$ rectangles
has a unique rectangle containing the northwest corner.  Extending its
southeast corner to the antidiagonal splits the remaining cells into two
smaller staircase regions; empty regions are allowed.  Recursing produces a
full ordered binary tree with one internal node per rectangle.  Conversely,
place the root rectangle across the corresponding antidiagonal interval and
recurse on the two subregions.  The split coordinates are integral and
bounded, so this is a primitive-recursive inverse pair.  It proves the
staircase-tiling interpretation stated on Wikipedia.

MathWorld also lists admissible diagonals in positive integral frieze patterns
and abstract flexagon configurations.  Their published correspondences pass
through triangulations or ordered pairs of ``pats.''  The triangulation side
has been proved above.  A complete formal proof for those specialized object
formats would require fixing their exact source conventions; this package
records them as formalization targets rather than silently changing the
definitions.  See the coverage ledger in Section~\ref{sec:coverage}.

\section{Formal series, boundary sums, special functions, and integrals}
\label{sec:series}

\subsection{The ordinary generating function}

Let $C(x)=\sum_{n\ge0}\Cat_nx^n$ as a formal series.  Segner's recurrence is
exactly
\[
 C(x)=1+xC(x)^2.                                                  \tag{4.1}\label{eq:Cquadratic}
\]
Multiplying the proposed closed form by its conjugate shows that the unique
solution with constant term $1$ is
\[
 C(x)=\frac{1-\sqrt{1-4x}}{2x}
     =\frac{2}{1+\sqrt{1-4x}}.                                   \tag{4.2}\label{eq:Cgf}
\]
Inside \ipra, Equation~\eqref{eq:Cgf} means the polynomial relation
Equation~\eqref{eq:Cquadratic}, the initial coefficient $1$, and the recursive
coefficient definition of the square root.  No choice between two completed
infinite series is needed: coefficient uniqueness performs the choice.

The elementary bound
\(\Cat_n\le\binom{2n}{n}\le 4^n\) proves absolute convergence for
$\abs{x}<1/4$.  The boundary is also absolutely convergent, by the exact tail
identity below.  Therefore the analytic equality in Equation~\eqref{eq:Cgf}
holds for complex $x$ with $\abs{x}\le1/4$, with the square root branch having
value $1$ at zero.

\subsection{An exact telescoping boundary identity}

Put $a_n=\Cat_n/4^n$.  The ratio recurrence gives
\[
 \frac{a_{n+1}}{a_n}=\frac{2n+1}{2n+4}.
\]
Hence, with $b_n=2(n+1)a_n$,
\[
 b_n-b_{n+1}=a_n.                                                 \tag{4.3}\label{eq:telescoping-term}
\]
This yields more than the infinite sum printed on the source pages.

\begin{theorem}[Exact partial sum and primitive-recursive tail]
\label{thm:boundary-tail}
For every $N\ge0$, writing $m=N+1$,
\begin{align}
 \sum_{n=0}^{N}\frac{\Cat_n}{4^n}
   &=2-\frac{2(m+1)\Cat_m}{4^m}
    =2-\frac{2\binom{2m}{m}}{4^m},                              \tag{4.4}\label{eq:partial-boundary}\\
 0\le 2-\sum_{n=0}^{N}\frac{\Cat_n}{4^n}
   &\le \frac{2}{\sqrt{m+1}}.                                  \tag{4.5}\label{eq:tail-bound}
\end{align}
Consequently
\[
 \sum_{n=0}^{\infty}\frac{\Cat_n}{4^n}=2,
 \qquad
 \sum_{n=1}^{\infty}\frac{\Cat_n}{4^n}=1.                      \tag{4.6}\label{eq:boundary-sums}
\]
For error at most $1/q$, it is enough to take $m+1\ge4q^2$.
\end{theorem}

\begin{proof}
Equation~\eqref{eq:partial-boundary} is the finite telescope of
Equation~\eqref{eq:telescoping-term}, followed by
$(m+1)\Cat_m=\binom{2m}{m}$.

For the bound, put $u_m=\binom{2m}{m}/4^m$.  We prove
$u_m^2\le1/(m+1)$ by induction.  It is true at $m=0$, and
\[
 \frac{u_{m+1}}{u_m}=\frac{2m+1}{2m+2}.
\]
The induction step is the integer polynomial inequality
\[
 (2m+1)^2(m+2)\le4(m+1)^3,
 \qquad
 4(m+1)^3-(2m+1)^2(m+2)=3m+2.
\]
Thus the tail $2u_m$ has square at most $4/(m+1)$.  If
$m+1\ge4q^2$, nonnegativity and square comparison give $2u_m\le1/q$.
This is already a denominator-cleared \ipra{} modulus; the square-root notation
in Equation~\eqref{eq:tail-bound} is only expository.
\end{proof}

\subsection{The random-walk interpretation}

A simple symmetric walk first reaches $-1$ at time $2n+1$ exactly when its
first $2n$ steps form a Dyck path after reversing the sign convention and the
last step is downward.  There are $\Cat_n$ such step strings, each of
probability $2^{-(2n+1)}$.  Therefore the probability of first capture by time
$2N+1$ is
\[
 \sum_{n=0}^{N}\frac{\Cat_n}{2^{2n+1}}
 =1-\frac{(N+2)\Cat_{N+1}}{4^{N+1}},                              \tag{4.7}
\]
with remaining probability at most $1/\sqrt{N+2}$.  The limit is $1$ by
Theorem~\ref{thm:boundary-tail}.  Thus the recurrence of the one-dimensional
walk is not imported as an external probabilistic theorem; it is supplied
with a finite exact probability and an explicit modulus.

\subsection{The exponential generating function}

Define
\[
 F(x)=\sum_{n\ge0}\Cat_n\frac{x^n}{n!}.
\]
The factorial majorant $\Cat_n\le4^n$ gives a primitive-recursive tail bound
for every bounded complex $x$, so $F$ is entire in the constructive sense.
Multiplying Equation~\eqref{eq:ratio-cleared} by $x^n/n!$ and comparing
coefficients gives the differential equation
\[
 xF''+(2-4x)F'-2F=0,
 \qquad F(0)=F'(0)=1.                                            \tag{4.8}\label{eq:Fode}
\]
Let
\[
 I_0(2x)=\sum_{j\ge0}\frac{x^{2j}}{(j!)^2},
 \qquad
 I_1(2x)=\sum_{j\ge0}\frac{x^{2j+1}}{j!(j+1)!}.                 \tag{4.9}
\]
Termwise finite coefficient arithmetic shows that
$G(x)=e^{2x}(I_0(2x)-I_1(2x))$ satisfies Equation~\eqref{eq:Fode} and the same
initial coefficients.  The coefficient recurrence from the ODE is unique,
so
\[
 \sum_{n\ge0}\Cat_n\frac{x^n}{n!}
 =e^{2x}\bigl(I_0(2x)-I_1(2x)\bigr).                             \tag{4.10}\label{eq:egf}
\]
This is a formal coefficient proof first; convergence upgrades it to every
complex $x$.  For $\abs{x}\le A$, once $j\ge2A$ the ratio of successive terms
in either Bessel series is at most $1/4$, so four thirds of the first omitted
term is a primitive-recursive uniform tail bound.

The second Bessel identity is even simpler.  Equation~\eqref{eq:factorial}
gives
\[
 \frac{\Cat_n}{(2n)!}=\frac1{n!(n+1)!},
\]
whence, coefficient by coefficient,
\[
 \sum_{n\ge0}\Cat_n\frac{x^{2n}}{(2n)!}
 =\frac{I_1(2x)}x,                                                \tag{4.11}\label{eq:even-egf}
\]
where the value at $x=0$ is the constant coefficient $1$.

\subsection{Integral and moment representations}

\begin{theorem}[Catalan beta integrals]
For $n\ge0$,
\[
 \Cat_n=\frac1{2\pi}\int_0^4 x^n\sqrt{\frac{4-x}{x}}\,dx
 =\frac{2}{\pi}4^n\int_{-1}^{1}t^{2n}\sqrt{1-t^2}\,dt.           \tag{4.12}\label{eq:integrals}
\]
\end{theorem}

\begin{proof}
In the first integral substitute $x=4u$.  The result is
\[
 \frac{2\cdot4^n}{\pi}
 B\left(n+\frac12,\frac32\right)
 =\frac{2\cdot4^n}{\pi}
 \frac{\Gamma(n+1/2)\Gamma(3/2)}{\Gamma(n+2)}.
\]
Use $\Gamma(3/2)=\sqrt\pi/2$ and the finite half-integer gamma calculation in
Equation~\eqref{eq:gamma}.  This gives $\Cat_n$.  For the second formula, use
$x=4t^2$ on $[0,1]$ and symmetry.

For a direct constructive implementation, use the second integral.  The
integrand is continuous on the compact interval after defining its endpoint
values to be zero.  For $f_n(t)=t^{2n}\sqrt{1-t^2}$ and
$\abs{s-t}\le\delta$ on $[-1,1]$,
\[
 \abs{f_n(s)-f_n(t)}\le 2n\delta+\sqrt{2\delta}.
\]
This follows from the telescoping bound for $s^{2n}-t^{2n}$ and
$\abs{\sqrt a-\sqrt b}\le\sqrt{\abs{a-b}}$.  Hence a uniform mesh of
width at most
\[
 \delta=\frac{1}{2048q^2(n+1)^2 16^n}
\]
gives upper and lower rational Riemann sums, after rational square-root
bisection, whose contribution to the resulting Catalan interval has width
below $1/(4q)$.  Approximate $2/\pi$ by its Machin-series interval finely
enough that the product contributes below $1/(2q)$; the combined width is
then below $1/q$.  Alternatively, repeated integration by parts reduces the integral to the
Wallis recurrence and Equation~\eqref{eq:factorial}, so only rational
endpoint data remain.
\end{proof}

The first density
\[
 d\mu(x)=\frac1{2\pi}\sqrt{\frac{4-x}{x}}\,\ind{0<x<4}\,dx       \tag{4.13}\label{eq:moment-density}
\]
has mass $\Cat_0=1$ and moments $\Cat_n$.  Thus the Catalan sequence is a
compactly supported Stieltjes moment sequence on $[0,4]$.  To state the
Hausdorff version precisely, put $x=4y$ and let
\[
 d\nu(y)=\frac2\pi\sqrt{\frac{1-y}{y}}\,\ind{0<y<1}\,dy.
\]
Then $\nu$ is a probability measure on $[0,1]$ and
\[
 \int_0^1 y^n\,d\nu(y)=\frac{\Cat_n}{4^n}.
\]
Hence the normalized Catalan sequence $(\Cat_n/4^n)$ is a Hausdorff moment
sequence, which is the precise compact-interval formulation behind the source
page's ``version of the Hausdorff moment problem.''  These formulas also give
positivity of every finite Hankel quadratic form, although
Section~\ref{sec:hankel} proves the exact determinants without analysis.

The boundary sum follows from the integral by geometric summation under the
integral, but Theorem~\ref{thm:boundary-tail} is stronger for \ipra: it avoids
interchanging an integral and an infinite series and supplies a small exact
tail.

\section{Divisibility, parity, prime Catalans, and final digits}
\label{sec:number-theory}

\subsection{Legendre and Kummer in bounded arithmetic}

For a prime $p$, let $\vp p(m)$ be the largest $e\le m$ such that $p^e$ divides
$m$; bounded search makes this primitive recursive.  Let $s_p(m)$ be the sum
of the base-$p$ digits of $m$, obtained by repeated quotient and remainder.

\begin{lemma}[Legendre]
\label{lem:legendre}
For every prime $p$ and $m\ge0$,
\[
 \vp p(m!)=\sum_{j=1}^{m}\floor{\frac{m}{p^j}}
          =\frac{m-s_p(m)}{p-1}.                                 \tag{5.1}\label{eq:legendre}
\]
\end{lemma}

\begin{proof}
In the product $m!$, every multiple of $p$ contributes once, every multiple of
$p^2$ contributes one additional time, and so on.  The sum may stop at $m$
because the later floors are zero.  This is an equality of bounded counts.
For the digit form, write $m=\sum d_ip^i$ and interchange the two finite sums:
\[
 \sum_{j\ge1}\floor{m/p^j}
 =\sum_i d_i(1+p+\cdots+p^{i-1})
 =\frac{m-\sum_i d_i}{p-1}.
\]
The internal statement is multiplied by $p-1$.
\end{proof}

The number of carries when adding $a$ and $b$ in base $p$ is
\((s_p(a)+s_p(b)-s_p(a+b))/(p-1)\).  Applying this with $a=b=n$ and
Equation~\eqref{eq:legendre} is the special Kummer theorem needed here.

\begin{theorem}[Prime valuation of a Catalan number]
\label{thm:valuation}
For every prime $p$,
\[
 \vp p(\Cat_n)
 =\frac{2s_p(n)-s_p(2n)}{p-1}-\vp p(n+1).                        \tag{5.2}\label{eq:valuation}
\]
Equivalently, it is the number of carries in the addition $n+n$ in base $p$,
minus the number of trailing zero digits of $n+1$.
\end{theorem}

\begin{proof}
Take $p$-adic valuations of
$(n+1)!n!\Cat_n=(2n)!$ and apply Lemma~\ref{lem:legendre}.  The quotient in
Equation~\eqref{eq:valuation} is exact because the digit-sum difference is a
multiple of $p-1$.  The carry interpretation follows by the usual digit
recurrence, which is a bounded induction over the digit positions.
\end{proof}

\subsection{Parity}

Since $s_2(2n)=s_2(n)$,
\[
 \vp 2(\Cat_n)=s_2(n)-\vp 2(n+1).                                \tag{5.3}
\]
Adding one changes the terminal block of $1$ bits to zeros, so
\[
 s_2(n+1)=s_2(n)+1-\vp 2(n+1).
\]
Therefore
\[
 \boxed{\ \vp 2(\Cat_n)=s_2(n+1)-1\ }.                          \tag{5.4}\label{eq:v2}
\]
It follows constructively that $\Cat_n$ is odd exactly when $n+1$ has one
binary $1$, that is,
\[
 \Cat_n\text{ is odd}\quad\Longleftrightarrow\quad n=2^k-1       \tag{5.5}\label{eq:odd}
\]
for some $k\le n+1$.  The witness $k$ is found by bounded search through
powers of two.

\subsection{The odd-prime digit automaton}

The verbal odd-prime rule on Wikipedia is most precise as a two-state scan.
Let $N=n+1=\sum_{i=0}^{L}d_ip^i$, with digits listed least significant first,
and let $t$ be the least index with $d_t\ne0$.  The $t$ initial zero digits
produce exactly $t$ carries when $N-1$ is doubled, and these are cancelled by
$\vp p(N)=t$.  Define carry bits $c_i$ for the remaining positions as follows:
\[
 c_t=
 \begin{cases}
  1,&t=0\text{ and }d_0\ge(p+3)/2,\\
  1,&t>0\text{ and }d_t\ge(p+1)/2,\\
  0,&\text{otherwise},
 \end{cases}                                                     \tag{5.6}\label{eq:carry-start}
\]
and, for $i>t$,
\[
 c_i=1\quad\Longleftrightarrow\quad
 d_i\ge\frac{p+1}{2}
 \ \text{or}\ 
 \left(d_i=\frac{p-1}{2}\ \text{and}\ c_{i-1}=1\right).        \tag{5.7}\label{eq:carry-step}
\]
Then
\[
 \vp p(\Cat_n)=\sum_{i=t}^{L}c_i.                               \tag{5.8}\label{eq:digit-automaton}
\]

\begin{proof}
The digits of $n=N-1$ are $p-1$ below $t$, $d_t-1$ at $t$, and $d_i$ above
$t$.  Doubling a digit $e_i$ with incoming carry $c$ produces an outgoing
carry exactly when $2e_i+c\ge p$.  For $i<t$ this always happens and is
cancelled by the subtraction of $t$ in Equation~\eqref{eq:valuation}.  At
$i=t$ the two cases in Equation~\eqref{eq:carry-start} distinguish whether there was
an incoming borrow-created carry.  Above $t$, the threshold is
$(p+1)/2$ without an incoming carry and $(p-1)/2$ with one, giving
Equation~\eqref{eq:carry-step}.  Bounded induction proves Equation~\eqref{eq:digit-automaton}.
\end{proof}

Read from the most significant end, Equations~\eqref{eq:carry-start}--\eqref{eq:carry-step}
are exactly the source-page rule: count digits above $(p+1)/2$; count a digit
equal to $(p+1)/2$ except at the least-significant starting position; and
propagate a count through a digit $(p-1)/2$ when the adjacent lower digit
produced a carry.

\subsection{Odd Catalans ending in five}

For an odd Catalan number, $n=2^k-1$ and $N=n+1=2^k$ is not divisible by $5$.
The automaton for $p=5$ begins at the least significant digit.  It produces no
carry precisely when that digit is at most $3$ and every higher base-five
digit is at most $2$.  Thus
\[
 5\nmid\Cat_{2^k-1}
 \Longleftrightarrow
 \begin{array}{l}
 \text{the base-five expansion of }2^k\text{ uses only }0,1,2,\\
 \text{except that its least significant digit may be }3.
 \end{array}                                                     \tag{5.9}\label{eq:base5-N}
\]
Subtracting one shows the equivalent MathWorld formulation:
\[
 5\nmid\Cat_{2^k-1}
 \Longleftrightarrow
 \text{the base-five expansion of }2^k-1\text{ uses only }0,1,2. \tag{5.10}\label{eq:base5-n}
\]
An odd integer ends in $5$ exactly when it is divisible by $5$, so these are
also the exact final-digit criteria.  The finite lists of exceptions reported
by the web pages are computational records; the supplied verifier checks a
configurable initial range from Equation~\eqref{eq:base5-N}, rather than
promoting a search range to a universal theorem.

\subsection{The only prime Catalan numbers}

\begin{lemma}
Every prime divisor of $\Cat_n$ is less than $2n$.
\end{lemma}

\begin{proof}
The factorial equation shows that a prime divisor of $\Cat_n$ divides
$(2n)!$, hence is at most $2n$.  Equality would make $2n$ prime.  For $n\ge2$
this is an even number greater than $2$, and for $n=1$ the Catalan number is
$1$.  Thus the inequality is strict.
\end{proof}

\begin{lemma}
For every $n\ge4$, $\Cat_n>2n-1$.
\end{lemma}

\begin{proof}
The base case is $\Cat_4=14>7$.  From
Equation~\eqref{eq:ratio-cleared}, for $n\ge1$,
\[
 \Cat_{n+1}\ge2\Cat_n,
\]
because $2(2n+1)\ge2(n+2)$.  If $\Cat_n>2n-1$, then
$\Cat_{n+1}>4n-2>2n+1=2(n+1)-1$ for $n\ge2$.
\end{proof}

\begin{theorem}[Catalan primes]
The only prime Catalan numbers are
\[
 \Cat_2=2\quad\text{and}\quad\Cat_3=5.                          \tag{5.11}
\]
\end{theorem}

\begin{proof}
The small values are direct.  If $n\ge4$ and $\Cat_n$ were prime, it would be
one of its own prime divisors and hence $\Cat_n<2n$, while the preceding lemma
gives the integer inequality $\Cat_n\ge2n$.  This is a contradiction.
\end{proof}

\section{Hankel determinants and the moment problem}
\label{sec:hankel}

The source page states three determinant formulas:
\[
 \det[\Cat_{i+j}]_{0\le i,j<n}=1,
 \qquad
 \det[\Cat_{i+j+1}]_{0\le i,j<n}=1,                              \tag{6.1}\label{eq:hankel01}
\]
and
\[
 \det[\Cat_{i+j+2}]_{0\le i,j<n}=n+1.                           \tag{6.2}\label{eq:hankel2}
\]
They have a particularly economical constructive proof.

\subsection{The ballot triangles}

For $0\le k\le r$, define
\[
 E_{r,k}=\binom{2r}{r-k}-\binom{2r}{r-k-1},                      \tag{6.3}\label{eq:even-ballot}
\]
and
\[
 O_{r,k}=\binom{2r+1}{r-k}-\binom{2r+1}{r-k-1}.                  \tag{6.4}\label{eq:odd-ballot}
\]
Outside the displayed range put the value zero.  Reflection proves that
$E_{r,k}$ counts nonnegative paths of $2r$ unit up/down steps ending at height
$2k$, and $O_{r,k}$ counts such paths of length $2r+1$ ending at height
$2k+1$.  Both arrays are lower triangular in $(r,k)$ and have diagonal entry
$1$.

Cut a Dyck path of semilength $r+s$ after $2r$ steps.  Its height there is
$2k$ for a unique $k$; the prefix is counted by $E_{r,k}$ and the reversed
suffix by $E_{s,k}$.  Therefore
\[
 \Cat_{r+s}=\sum_{k\ge0}E_{r,k}E_{s,k}.                          \tag{6.5}\label{eq:even-gram}
\]
Cutting a Dyck path of semilength $r+s+1$ after $2r+1$ steps similarly gives
\[
 \Cat_{r+s+1}=\sum_{k\ge0}O_{r,k}O_{s,k}.                        \tag{6.6}\label{eq:odd-gram}
\]
The cut, reversal, and gluing operations are primitive-recursive inverse maps,
so these are finite combinatorial identities, not analytic moment arguments.

\subsection{The two unit determinants}

Let $E^{(n)}=(E_{r,k})_{0\le r,k<n}$ and similarly $O^{(n)}$.  Equations
\eqref{eq:even-gram}--\eqref{eq:odd-gram} say
\[
 [\Cat_{i+j}]=E^{(n)}(E^{(n)})^{\mathsf T},
 \qquad
 [\Cat_{i+j+1}]=O^{(n)}(O^{(n)})^{\mathsf T}.                    \tag{6.7}
\]
Both factors are unit lower triangular.  Determinant multiplicativity and
transpose invariance now give Equation~\eqref{eq:hankel01}.  In a minimal
formalization, multiplicativity is needed only for triangular matrices and
their transposes; it can also be proved by bounded Gaussian elimination using
integer row operations.

\subsection{The determinant starting at \(\Cat_2\)}

Let $M^{(n)}$ be the $n$ by $(n+1)$ matrix
\[
 M^{(n)}_{r,k}=E_{r+1,k},
 \qquad 0\le r<n,
 \quad 0\le k\le n.
\]
Cutting after $2r+2$ steps gives
\[
 [\Cat_{r+s+2}]_{0\le r,s<n}=M^{(n)}(M^{(n)})^{\mathsf T}.       \tag{6.8}\label{eq:Mgram}
\]
Cauchy--Binet expresses its determinant as the sum of the squares of the
$n+1$ maximal minors of $M^{(n)}$.  Every one of those minors is $1$.

\begin{lemma}[Maximal minors of the shifted ballot matrix]
\label{lem:maximal-minors}
Deleting any one column of $M^{(n)}$ leaves an $n$ by $n$ matrix of determinant
$1$.
\end{lemma}

\begin{proof}
Deleting column $0$ leaves columns $1,\ldots,n$.  The result is unit lower
triangular because $E_{r+1,k}=0$ for $k>r+1$ and
$E_{r+1,r+1}=1$, so this minor is $1$.

For every $r\ge1$,
\[
 \sum_{k=0}^{r}(-1)^k E_{r,k}=0.                                 \tag{6.9}\label{eq:alternating-ballot}
\]
To verify this, substitute Equation~\eqref{eq:even-ballot}, shift the second
finite sum, and use the symmetry of the binomial coefficients in the expansion
of $(1-1)^{2r}$.  Hence the alternating vector
$(1,-1,1,\ldots,(-1)^n)^{\mathsf T}$ spans the kernel of $M^{(n)}$; the rank is
$n$ because the minor just exhibited is nonzero.

The signed cofactor vector
\[
 \bigl((-1)^j\det M^{(n)}_{\widehat j}\bigr)_{j=0}^{n}
\]
is also in that one-dimensional kernel.  Comparing its first coordinate with
the known unit minor shows it equals the alternating vector.  Thus every
unsigned maximal minor is $1$.  This argument uses only the Laplace expansion
and cancellation in $\Z$, all available as finite primitive-recursive
operations.
\end{proof}

Applying Cauchy--Binet to Equation~\eqref{eq:Mgram} now yields
\[
 \det[\Cat_{i+j+2}]_{0\le i,j<n}
 =\sum_{j=0}^{n}1^2=n+1,
\]
which is Equation~\eqref{eq:hankel2}.

\subsection{The two determinant conditions characterize the sequence}

Suppose a sequence $a_0,a_1,\ldots$ satisfies, for every positive $n$,
\[
 \det[a_{i+j}]_{0\le i,j<n}=1,
 \qquad
 \det[a_{i+j+1}]_{0\le i,j<n}=1.                                \tag{6.10}\label{eq:characterizing-Hankel}
\]
The $1$ by $1$ cases give $a_0=a_1=1$.  Assume $a_0,\ldots,a_{2r-1}$ have
already been determined.  In the first determinant of size $r+1$, the newest
term $a_{2r}$ occurs only in the southeast corner.  Its cofactor is the
preceding unshifted determinant, equal to $1$.  Therefore the equation
``determinant $=1$'' is a linear equation with coefficient $1$ and uniquely
determines $a_{2r}$.  The shifted determinant of size $r+1$ similarly
determines $a_{2r+1}$, with cofactor $1$.

Quantifier-free induction therefore proves uniqueness.  Since the Catalan
sequence satisfies Equation~\eqref{eq:characterizing-Hankel}, the two
conditions together characterize it, exactly as stated on Wikipedia.

\subsection{Positive Hankel forms and compact moment representations}

The Gram factorizations give a finite integral-free positivity theorem.  For
any integer vector $z=(z_0,\ldots,z_{n-1})$,
\[
 \sum_{i,j<n}z_iz_j\Cat_{i+j}
 =\sum_{k<n}\left(\sum_{i<n}z_iE_{i,k}\right)^2\ge0.             \tag{6.11}
\]
The determinant $1$ shows positive definiteness over $\Q$.  The density in Equation~\eqref{eq:moment-density} gives the compatible
compactly supported Stieltjes representation on $[0,4]$:
\[
 \sum_{i,j}z_iz_j\Cat_{i+j}
 =\int_0^4\left(\sum_i z_ix^i\right)^2d\mu(x).
\]
After $x=4y$, the normalized matrix
$[\Cat_{i+j}/4^{i+j}]$ has a Hausdorff representation on $[0,1]$.
Thus the finite combinatorial factorization and the constructive moment code
prove the same positivity by independent routes.

\section{Ballot, convolution, rational, Fuss, and super-Catalan extensions}
\label{sec:generalizations}

\subsection{Bertrand's ballot theorem at excess one}

Consider words with $n+1$ votes $A$ and $n$ votes $B$.  Give $A$ weight $+1$
and $B$ weight $-1$.  The cycle lemma says that a word with total weight
$r>0$ has exactly $r$ distinguished cyclic starting positions whose nonempty
partial sums are positive, counted with positions even when the underlying
word is periodic.  For total weight one, every pointed word has exactly one
good start.  Double counting pointed words gives
\[
 \frac1{2n+1}\binom{2n+1}{n}=\Cat_n.                             \tag{7.1}
\]
Removing the first $A$ from a positive word gives a Dyck word.  This proves the
ballot interpretation and Equation~\eqref{eq:odd-binomial} simultaneously.

For \ipra, the cycle lemma can be proved by the deletion algorithm mentioned
on Wikipedia: place the letters on a circle, repeatedly delete adjacent
$AB$ pairs, and record the surviving $A$ positions.  The deletion relation is
bounded, a fixed least-pair strategy makes it primitive recursive, and an
induction on word length proves that the survivors are exactly the good
starts.  For total weight $r$, exactly $r$ survive.

\subsection{The \(k\)-fold Catalan convolution}

A forest of $k$ ordered full binary trees with a total of $n$ internal nodes is
counted in two ways.  Splitting by the internal-node counts of the components
gives
\[
 \sum_{i_1+\cdots+i_k=n}\Cat_{i_1}\cdots\Cat_{i_k}.              \tag{7.2}
\]
Encode the forest in Polish order, writing an internal node with weight $+1$
and a leaf with weight $-1$.  There are $n$ internal nodes and $n+k$ leaves,
so the word length is $2n+k$ and its total weight is $-k$.  The Raney form of
the cycle lemma says that, among the $2n+k$ pointed rotations of every such
word, exactly $k$ have all proper prefix sums greater than $-k$ in the order
needed to cut the word uniquely into $k$ trees.  Hence
\[
 \boxed{\displaystyle
 \sum_{i_1+\cdots+i_k=n}\Cat_{i_1}\cdots\Cat_{i_k}
  =\frac{k}{2n+k}\binom{2n+k}{n}.}                               \tag{7.3}\label{eq:kfold}
\]
The internal theorem is
$(2n+k)$ times the left side equals $k\binom{2n+k}{n}$.
The pointed-word proof avoids any assumption that cyclic orbits have full
size; periodic words cause no difficulty because the distinguished position is
part of the code.

Coefficientwise, Equation~\eqref{eq:kfold} is also Lagrange inversion for
$C(x)^k$.  The forest proof is the smaller \ipra{} dependency.

\subsection{Rational Catalan numbers}

Let $r,s$ be coprime positive integers.  Give a right step weight $s$ and an up
step weight $-r$.  A path from $(0,0)$ to $(r,s)$ has total weight zero.  Mark
one step and perform the standard cycle-lemma normalization at the least
partial sum.  Coprimality implies that the $r+s$ cyclic partial-sum ranks are
distinct modulo $r+s$, so exactly one marked rotation stays on the required
side of the line $ry=sx$.  Double counting gives
\[
 \operatorname{RatCat}(r,s)
 =\frac1{r+s}\binom{r+s}{r}.                                    \tag{7.4}\label{eq:rational-cat}
\]
The proof is constructive: compute all $r+s$ partial sums, choose the least
rank with the fixed tie convention, rotate, and invert by remembering the
mark.  Coprimality is used through Bezout data obtained by the primitive-
recursive extended Euclidean algorithm.  Equation~\eqref{eq:rational-cat}
is represented after multiplication by $r+s$.

The ordinary Catalan case is $(r,s)=(n,n+1)$ after the strict/weak boundary
translation; Equation~\eqref{eq:odd-binomial} is the resulting formula.

\subsection{Fuss--Catalan numbers}

For $p\ge2$ and $k\ge0$, define
\[
 \Cat^{(p)}_0=1,
 \qquad
 \Cat^{(p)}_k
 =\frac1{(p-1)k+1}\binom{pk}{k}
 =\frac1k\binom{pk}{k-1}\quad(k\ge1).                           \tag{7.5}\label{eq:fuss}
\]
The two fractions are equal by adjacent-binomial arithmetic.

A full ordered $p$-ary tree with $k$ internal nodes has $(p-1)k+1$ leaves.
Its Polish word therefore has $k$ symbols of weight $p-1$ and
$(p-1)k+1$ symbols of weight $-1$, total length $pk+1$ and total weight $-1$.
The cycle lemma gives exactly one valid pointed rotation.  Thus the number of
trees is
\[
 \frac1{pk+1}\binom{pk+1}{k}
 =\frac1{(p-1)k+1}\binom{pk}{k},
\]
proving Equation~\eqref{eq:fuss}.  Recursive parsing of the Polish word is the
primitive-recursive inverse.

The same trees are syntax trees for $k$ applications of a fixed $p$-ary
operator.  Rooting a dissection at a boundary edge and taking the incident
$(p+1)$-gon gives a $p$-tuple of smaller dissections, so they also count
partitions of a convex $((p-1)k+2)$-gon into $k$ nonoverlapping
$(p+1)$-gons.  The Lukasiewicz height of the Polish word gives the $p$-good
path interpretation.  These are explicit recursive bijections, not merely
recurrence comparisons.

The formal series $T=1+xT^p$ has coefficients $\Cat^{(p)}_k$.  For every $q\ge1$, finite Lagrange inversion gives
\[
 [x^k]T^q=\frac{q}{pk+q}\binom{pk+q}{k},                         \tag{7.6}\label{eq:fuss-power}
\]
with the denominator cleared.  Forests of $q$ $p$-ary trees give the same
identity by the Raney lemma.

\subsection{The two-parameter path numbers on MathWorld}

The notation on MathWorld can be handled uniformly by formal reversion.  Fix
$p>1$ and $0\le q\le p-1$.  Let $X=X(y)$ be the unique formal series with
\[
 X=y(1+X)^p.                                                      \tag{7.7}\label{eq:Xreversion}
\]
Define
\[
 D_{p,q}(y)=\frac{X}{(1+X)^q}=\sum_{k\ge1}d_{p,q}(k)y^k.          \tag{7.8}\label{eq:Dpq}
\]
Finite Lagrange inversion gives
\[
 d_{p,q}(k)=\frac{p-q}{pk-q}\binom{pk-q}{k-1}.                  \tag{7.9}\label{eq:Dexplicit}
\]
For $q=0$ this is the Fuss--Catalan number in
Equation~\eqref{eq:fuss}.  The coefficient proof is a bounded sum: from
\([y^k]\Phi(X)=k^{-1}[t^{k-1}]\Phi'(t)(1+t)^{pk}\), take
$\Phi(t)=t(1+t)^{-q}$ and simplify by Pascal's identity.  Negative exponents
are avoided by multiplying by a sufficiently large power of $1+t$ before
coefficient extraction.

Substituting $y=x/(1+x)^p$ in Equation~\eqref{eq:Dpq} gives
\[
 D_{p,q}\left(\frac{x}{(1+x)^p}\right)=\frac{x}{(1+x)^q}.        \tag{7.10}\label{eq:Dsub}
\]
Taking the coefficient of $x^k$ after multiplication by $(1+x)^n$ proves the
generalized Jonah identity displayed by MathWorld:
\[
 \binom{n-q}{k-1}
 =\sum_{i=1}^{k}d_{p,q}(i)\binom{n-pi}{k-i}.                    \tag{7.11}\label{eq:jonah}
\]
The natural generalized-binomial convention applies if an upper entry is
negative; in the ordinary path range all terms are nonnegative and finite.

For $r>0$ with $q<p-r$, Equation~\eqref{eq:Dpq} also gives
\[
 D_{p,q}(y)=\frac{D_{p,p-r}(y)D_{p,q+r}(y)}{y}.                  \tag{7.12}
\]
Comparing the coefficient of $y^k$ yields the MathWorld recurrence
\[
 d_{p,q}(k)=
 \sum_{\substack{i,j>0\\i+j=k+1}}
 d_{p,p-r}(i)d_{p,q+r}(j).                                      \tag{7.13}\label{eq:Dconvolution}
\]
Equations~\eqref{eq:Dexplicit}, \eqref{eq:jonah}, and
\eqref{eq:Dconvolution} therefore have a single coefficientwise \ipra{} proof.
The path interpretation follows by decomposing at the first contact with the
appropriate $p$-slope boundary; the formal proof is independent of drawing
conventions.

\subsection{Super-Catalan integers}

The two-parameter super-Catalan numbers on Wikipedia are
\[
 S(m,n)=\frac{(2m)!(2n)!}{m!n!(m+n)!}.                           \tag{7.14}\label{eq:supercat}
\]
Their integrality is not obvious from the quotient, but it has a short
valuation proof suitable for \ipra.  For every prime power $a=p^j$,
\[
 \floor{\frac{2m}{a}}+\floor{\frac{2n}{a}}
 \ge
 \floor{\frac m a}+\floor{\frac n a}+\floor{\frac{m+n}{a}}.     \tag{7.15}\label{eq:super-floor}
\]
After removing integer parts, this is
$\floor{2u}+\floor{2v}\ge\floor{u+v}$ for $0\le u,v<1$, verified by
the four cases $u<1/2$ or $u\ge1/2$, and similarly for $v$.  Summing
Equation~\eqref{eq:super-floor} over $j$ with Legendre's formula shows
$\vp p(S(m,n))\ge0$ for every prime $p$, hence the denominator divides the
numerator.  In bounded arithmetic, a least-prime-divisor induction supplies
the final divisibility conclusion.

Direct cancellation gives
\[
 S(1,n)=2\Cat_n,
 \qquad
 S(n,n)=\binom{2n}{n}.                                           \tag{7.16}
\]
The source correctly notes that a uniform combinatorial interpretation for
all $(m,n)$ is open.  This report proves integrality and the displayed special
cases but does not manufacture a nonexistent general bijection.

\subsection{Coxeter Catalan statements: type \(A\) and two source corrections}

For the standard generalized Catalan number of the rank-$n$ root system
$A_n$ \cite{FominReading}, identify a positive root with an interval $[i,j]$.  Order ideals in the
triangular interval poset are in primitive-recursive bijection with Dyck paths:
trace the northeast boundary of the ideal.  Taking maximal elements and taking
downward closure are inverse bounded maps between order ideals and antichains.
Consequently
\[
              \operatorname{Cat}(A_n)=\Cat_{n+1}.               \tag{7.17}
\]
Thus, in the standard rank convention, the ordinary number $\Cat_n$
corresponds to $A_{n-1}$, not $A_n$.

Two sentences on the requested Wikipedia revision need correction.
First, root-poset antichains are not equinumerous with fully commutative group
elements in every finite crystallographic type \cite{StembridgeFC}.  The smallest counterexample
is $B_2=I_2(4)$.  Its four positive roots have the two simple roots incomparable below $\alpha_1+\alpha_2$, which is below the highest root; their antichains are
\[
 \varnothing,\quad \{\alpha_1\},\quad \{\alpha_2\},\quad
 \{\alpha_1+\alpha_2\},\quad \{\alpha_1+2\alpha_2\},\quad
 \{\alpha_1,\alpha_2\},
\]
so there are $6$.  The dihedral Coxeter group has $8$ elements, and every
one except the longest element is fully commutative, so there are $7$.
The two counts do agree in type $A$: fully commutative permutations are the
$321$-avoiding permutations already counted in Section~\ref{sec:combinatorics}.

Second, the unweighted assertion that a standard generalized Catalan number
is the sum over all maximal proper subdiagrams is false.  For $A_2$ the left
side is $\Cat_3=5$, whereas the two deleted diagrams are $A_1$, each with
Catalan number $\Cat_2=2$, giving $4$.  In type $A$ the correct weighted
maximal-parabolic recurrence is
\[
 2n\,\operatorname{Cat}(A_n)
  =(n+3)\sum_{i=1}^{n}
      \operatorname{Cat}(A_{i-1})\operatorname{Cat}(A_{n-i}).  \tag{7.18}\label{eq:coxeter-A-recurrence}
\]
Indeed, with $\operatorname{Cat}(A_j)=\Cat_{j+1}$ and
$\operatorname{Cat}(A_0)=1$, the sum is
\[
 S_n=\sum_{i=1}^{n}\Cat_i\Cat_{n+1-i}
     =\Cat_{n+2}-2\Cat_{n+1}
     =\frac{2n}{n+3}\Cat_{n+1},
\]
where the middle equality is Segner and the last is the ratio recurrence.
Clearing $n+3$ gives Equation~\eqref{eq:coxeter-A-recurrence} as a finite
\ipra{} identity.  The unweighted Segner sum does hold for the shifted
``positive'' type-$A$ Catalan invariant $\operatorname{Cat}^{+}(A_n)=\Cat_n$;
this likely explains the conflation on the source page.

The general root-poset theorem and the correctly normalized parabolic
recurrence for every finite crystallographic type require the classification
of finite root systems (or a uniform Coxeter-theoretic development).  They are
recorded as formalization targets, not asserted as consequences of the false
unweighted sentence.

\section{Effective asymptotics and decimal digits}
\label{sec:asymptotics}

The expression $u_n\sim v_n$ is not taken as primitive notation.  We prove
rational inequalities with explicit moduli for the Cauchy codes of $\pi$ and
square roots.

\subsection{A Wallis bound and the leading equivalent}

Let
\[
 q_n=\frac1{4^n}\binom{2n}{n}.
\]
The finite Wallis inequalities give, for $n\ge1$,
\[
 \frac1{\sqrt{\pi(n+1/2)}}<q_n<\frac1{\sqrt{\pi n}}.             \tag{8.1}\label{eq:wallis-central}
\]
For completeness, one constructive derivation starts with
$I_m=\int_0^{\pi/2}\sin^m t\,dt$.  Integration by parts gives
$I_m=(m-1)I_{m-2}/m$, while
$I_{2n+1}<I_{2n}<I_{2n-1}$ follows pointwise on $(0,\pi/2)$.  Substitute the
finite product formulas for $I_{2n}$ and $I_{2n+1}$, square, and clear all
positive factors.  Refining the elementary comparison with
$\sin t$ on rational polygonal meshes gives Equation~\eqref{eq:wallis-central}
and a primitive-recursive Cauchy code for $\pi$.  No impredicative completeness
principle enters the proof.

Define the normalized ratio
\[
 R_n=\frac{\sqrt\pi\,n^{3/2}\Cat_n}{4^n}
     =\frac n{n+1}\sqrt{\pi n}\,q_n.                            \tag{8.2}
\]
Equation~\eqref{eq:wallis-central} implies
\[
 \frac n{n+1}\sqrt{\frac n{n+1/2}}<R_n<\frac n{n+1}<1.          \tag{8.3}\label{eq:Rbounds}
\]
The rational inequalities
\[
 \frac n{n+1}\ge1-\frac1n,
 \qquad
 \frac1{\sqrt{1+x}}\ge1-\frac x2\quad(x\ge0)
\]
then give
\[
 0<1-R_n<\frac{5}{4n}<\frac2n.                                  \tag{8.4}\label{eq:leading-error}
\]
The square-root inequality is proved by squaring after checking both sides are
nonnegative.  Therefore a convenient primitive-recursive modulus is
\[
 n\ge2q\quad\Longrightarrow\quad \abs{R_n-1}\le\frac1q.         \tag{8.5}\label{eq:leading-modulus}
\]
This proves, in the exact constructive sense,
\[
 \boxed{\displaystyle
 \Cat_n\sim\frac{4^n}{\sqrt\pi\,n^{3/2}}.}                       \tag{8.6}
\]

\subsection{The first three displayed asymptotic terms}

The source page prints
\[
 \Cat_n\sim\frac{4^n}{\sqrt\pi}
 \left(n^{-3/2}-\frac98n^{-5/2}+\frac{145}{128}n^{-7/2}+\cdots\right).
\tag{8.7}\label{eq:printed-expansion}
\]
We now give an effective remainder.

\begin{lemma}[Two-term effective Stirling logarithm]
For $m\ge1$,
\[
 \log(m!)=
 \left(m+\frac12\right)\log m-m+\frac12\log(2\pi)
 +\frac1{12m}-\frac1{360m^3}+\rho_m,                            \tag{8.8}
\]
where
\[
 0<\rho_m<\frac1{1260m^5}.                                      \tag{8.9}
\]
\end{lemma}

\begin{proof}
Apply Euler--Maclaurin to $\log x$ through the Bernoulli-$B_4$ term.  The next
derivative has fixed sign, so the integral remainder lies between zero and
the first omitted Bernoulli term.  At this fixed order the Bernoulli
polynomials, their maxima on $[0,1]$, and the integral comparison are rational
polynomial calculations.  For an object-language proof, replace the integral
by upper and lower rational Riemann sums; the derivative bound gives a mesh
modulus primitive recursive in $(m,q)$.  The endpoint constant is fixed by the
Wallis code for $\pi$.  This is the standard effective proof of Stirling, but
all quantifiers have explicit finite bounds.
\end{proof}

Subtracting twice the formula at $n$ from the formula at $2n$ gives
\[
 \log\left(\frac{\binom{2n}{n}\sqrt{\pi n}}{4^n}\right)
 =-\frac1{8n}+\frac1{192n^3}+\eta_n,
 \qquad \abs{\eta_n}\le\frac1{600n^5}.                          \tag{8.10}\label{eq:central-log}
\]
Put $u=1/n$.  Since $R_n$ is the central-binomial ratio multiplied by
$1/(1+u)$,
\[
 \log R_n=-\frac98u+\frac12u^2-\frac{21}{64}u^3+\delta_n,
 \qquad \abs{\delta_n}\le\frac{13}{50}u^4                     \tag{8.11}\label{eq:Rlog}
\]
for $n\ge8$.  Here the alternating expansion of $\log(1+u)$ has remainder at
most $u^4/4$, and Equation~\eqref{eq:central-log} supplies the smaller
$u^5/600$ term.

For $n\ge8$, Equation~\eqref{eq:Rlog} gives $\abs{\log R_n}<3/20$.
Here is an explicit rational audit of the Taylor constant.  Put
$z=\log R_n$, $a_1=-9/8$, and $a_2=1/2$.  Since $u\le1/8$,
Equation~\eqref{eq:Rlog} implies
\[
 |z|\le\frac65u,\qquad
 |z-a_1u|\le\frac{11}{20}u^2,\qquad
 |z+a_1u|\le\frac{12}{5}u.
\]
Also $e^{3/20}<7/6$, as follows by summing its positive Taylor series with a
geometric tail after the linear term.  Taylor's theorem therefore gives
\begin{align*}
 \left|e^z-\left(1+a_1u+
          \left(a_2+\frac{a_1^2}{2}\right)u^2\right)\right|
 &\le \left(\frac{21}{64}+\frac{13}{400}\right)u^3
   +\frac12\frac{11}{20}\frac{12}{5}u^3
   +\frac{7}{36}\left(\frac65\right)^3u^3\\
 &<2u^3.
\end{align*}
The three summands respectively bound the cubic-and-$\delta_n$ part of $z$,
the correction from $z^2-a_1^2u^2$, and the exponential remainder.  Hence
\[
 \boxed{\displaystyle
 \abs{R_n-
  \left(1-\frac9{8n}+\frac{145}{128n^2}\right)}
 \le\frac2{n^3}\qquad(n\ge8).}                                 \tag{8.12}\label{eq:three-term-bound}
\]
The coefficient calculation is
\[
 \frac12+\frac12\left(\frac98\right)^2=\frac{145}{128}.
\]
Collecting one further coefficient gives the next term
$-1155/(1024n^3)$.  It is not needed below; the report stops at
Equation~\eqref{eq:three-term-bound} because that bound exactly certifies all
terms printed by MathWorld.

Every transcendental operation in this argument is evaluated by a rational
interval.  For a requested accuracy $1/q$, choose the Taylor lengths by the
explicit alternating/geometric tail bounds, and choose $n$ from the displayed
polynomial inequalities.  Thus the proof has a primitive-recursive modulus
and is admissible in the constructive setting.

\subsection{Decimal digit ratios}

Let $d(N)$ be the number of decimal digits of $\Cat_N$, so
$d(N)=\floor{\log_{10}\Cat_N}+1$.  From Equation~\eqref{eq:Rbounds}, for
$N\ge2$,
\[
 \log_{10}\Cat_N
 =N\log_{10}4-\frac32\log_{10}N-\frac12\log_{10}\pi
  +\log_{10}R_N.                                                 \tag{8.13}
\]
Since $1/2<R_N<1$ for $N\ge3$, the last two terms have bounded absolute value.
Together with the one-unit floor error, a convenient explicit estimate is
\[
 \abs{\frac{d(N)}N-\log_{10}4}
 \le\frac{\tfrac32\log_{10}N+3}{N}\qquad(N\ge3).                \tag{8.14}\label{eq:digit-bound}
\]
In particular, at $N=10^m$ with $m\ge1$,
\[
 \abs{\frac{d(10^m)}{10^m}-\log_{10}4}
 \le\frac{\tfrac32m+3}{10^m},                                  \tag{8.15}\label{eq:digit-ratio}
\]
which proves the convergence claimed by MathWorld with a primitive-recursive
modulus.  The logarithm base ten is a rational Cauchy code obtained from the
series for $\log(1+x)$ after range reduction.  The finite digit values printed
on the source page are computational data; the supplied checker independently
certifies those within its configured feasible range, while
Equation~\eqref{eq:digit-ratio} proves the general limiting statement.

\section{Lean formalization blueprint}
\label{sec:formalization}

\subsection{The exact checked starting point}

The inspected repository snapshot is
\texttt{00bbe871268c453903499fb6afe07519e87d9631}.  The relevant namespace is
\begin{center}
\texttt{LeanProofs.IPRA.Combinatorics.BinomialIdentities}.
\end{center}
It imports the exact binomial development and proves, without
\texttt{sorry}, the following named results:
\begin{itemize}[leftmargin=2em]
\item \texttt{binomSumFormula}, \texttt{binom\_sum}: the row sum
      $\sum_{k\le n}\binom nk=2^n$;
\item \texttt{centralFormula}, \texttt{central\_binomial}: the central
      binomial inequality used by that module;
\item \texttt{catalanRatioFormula}, \texttt{catalan\_ratio}: the adjacent
      binomial identity;
\item \texttt{catalanFormula}, \texttt{catalan}: the universal closure of
      Equation~\eqref{eq:checked-catalan}.
\end{itemize}
The binomial implementation in
\texttt{Combinatorics/Binomial.lean} defines
\(\binom nk=n!\mathbin{\mathrm{quo}}(k!(n\monus k)!)\), proves exactness by a
bounded remainder sum, and derives Pascal's rule internally.  This is the
right foundation for the present project; redefining binomial coefficients in
ambient \texttt{Nat} would weaken the formal result.

A minimal reuse file begins as follows.
\begin{lstlisting}[language={},caption={Checked base import and Catalan term}]
import IPRADerivations.Combinatorics.BinomialIdentities

namespace LeanProofs.IPRA.Combinatorics.Catalan

open LeanProofs.ConciseIPRA LeanProofs.IPRA
open LeanProofs.IPRA.Combinatorics.Binomial
open LeanProofs.IPRA.Combinatorics.BinomialIdentities
open Syn
open scoped Fm
open ConciseIPRA.PTm (app zero succ)
open LeanProofs.IPRA.PTm (num monus)

abbrev cat {m : Nat} (n : PTm m) : PTm m :=
  monus (binom (n + n) n) (binom (n + n) (succ n))

-- Existing checked theorem:
#check LeanProofs.IPRA.Combinatorics.BinomialIdentities.catalan
\end{lstlisting}
The supplied \texttt{lean/CatalanProveItBlueprint.lean} records this import and
unproved target formulas without inserting axioms or placeholders.

\subsection{Proposed module graph}

\begin{longtable}{@{}p{0.27\textwidth}p{0.66\textwidth}@{}}
\toprule
\textbf{New module} & \textbf{Responsibility and dependencies}\\
\midrule
\endfirsthead
\toprule
\textbf{New module} & \textbf{Responsibility and dependencies}\\
\midrule
\endhead
\path{Catalan/Core.lean} & Define \texttt{cat} by binomial difference;
prove $n=0$ and instantiate the checked \texttt{catalan}; derive factorial,
odd-binomial, product, and ratio equations.  Depends only on
\texttt{BinomialIdentities}.\\
\path{Catalan/Dyck.lean} & Bit-word coding, prefix heights, Dyck Boolean,
reflection at the first negative prefix, first-return split, and Segner.
Depends on pairing, bounded search, and \texttt{Catalan/Core}.\\
\path{Catalan/ProofVariants.lean} & Chung--Feller cut-and-swap with its
backward-scan inverse, marked triangulation expansion/collapse, and the
balanced-string convolution-cancellation proof.  Depends on
\texttt{Catalan/Dyck} and \texttt{Catalan/Bijections}.\\
\path{Catalan/Bijections.lean} & Binary and plane trees, parentheses,
lattice paths, triangulations, matchings, restricted sequences, and tableaux.
Each class is a bounded code with a \texttt{Valid} formula and explicit
\texttt{encode}/\texttt{decode} inverse lemmas.\\
\path{Catalan/FiniteSums.lean} & Two-colored Motzkin, weighted convolution,
alternating transform, square sum, Narayana/hypergeometric identity, and
Riordan transform.  Depends on \texttt{Core/SumLemmas} and Vandermonde.\\
\path{Catalan/Permutations.lean} & Bounded permutation codes,
standardization, $231$ decomposition, stack sorting, finite RSK, and all
length-three patterns.\\
\path{Catalan/Digits.lean} & Base-$p$ digits, digit sum, carries, Legendre,
valuation formula, parity, base-five automaton, and prime Catalans.  Reuses
\texttt{AnalyticNumberTheory/Primes}.\\
\path{Catalan/Hankel.lean} & Ballot triangles, two Gram identities,
triangular determinants, rectangular Cauchy--Binet, maximal-minor lemma, and
uniqueness.  Reuses \texttt{Core/LinearAlgebra}.\\
\path{Catalan/FormalSeries.lean} & Finite coefficient arrays, Cauchy
product, substitution with zero constant term, $C=1+xC^2$, Bessel coefficient
identities, and Lagrange inversion at fixed coefficient.\\
\path{Catalan/Generalized.lean} & Cycle/Raney lemma, forests, rational and
Fuss--Catalan numbers, generalized Jonah identity, convolution, and
super-Catalan integrality.\\
\path{Catalan/CoxeterAudit.lean} & Type-$A$ interval-root posets, the weighted
maximal-parabolic recurrence, and the finite $A_2/B_2$ source counterexamples.
The all-types theorem remains a separate dependency-heavy target.\\
\path{Catalan/RealCodes.lean} & Rational interval codes for $\pi$, square
root, logarithm, exponential, Riemann sums, Wallis, Stirling remainder,
integrals, and asymptotic moduli.  This is the only genuinely new analytic
infrastructure.\\
\bottomrule
\end{longtable}

\subsection{Representative object-language formulas}

The ratio recurrence should be stored without division.  With variables
$0=n$, a target formula is
\begin{lstlisting}[language={}]
abbrev ratioFormula : Fm PTm 1 :=
  ((succ (succ (var 0))) * cat (succ (var 0)))
    ≐ ((num 2 * succ (var 0 + var 0)) * cat (var 0))
\end{lstlisting}
Here $\doteq$ (written as \texttt{U+2250} in source) is the object-language equality constructor used by the
imported ProveIt modules; the theorem target is the universal closure of this
formula in \texttt{Prf PRAAx}.

For Segner, define a two-variable term body
\[
 b(k,n)=\Cat_k\Cat_{n\monus k},
\]
then use the existing primitive-recursive bounded sum symbol
\texttt{sumF}:
\begin{lstlisting}[language={}]
abbrev segnerBody : PTm 2 := cat (var 0) * cat (monus (var 1) (var 0))
abbrev segnerRhs (n : PTm m) : PTm m :=
  ap2 (sumF segnerBody) (succ n) n
abbrev segnerFormula : Fm PTm 1 :=
  cat (succ (var 0)) ≐ segnerRhs (var 0)
\end{lstlisting}
The proof can either formalize first return or show that both sides satisfy the
checked ratio recurrence.  The bijective route is preferable because later
modules reuse the split maps.

A uniform finite-cardinality API should use
\begin{lstlisting}[language={}]
structure BoundedClass where
  bound : Nat -> Nat
  valid : Nat -> Nat -> Bool

def count (A : BoundedClass) (n : Nat) : Nat :=
  (List.range (A.bound n)).foldl
    (fun s z => s + if A.valid n z then 1 else 0) 0
\end{lstlisting}
only in the executable layer.  The object-language layer mirrors these
functions as primitive-recursive symbols and proves a generic bijection lemma
once.  Subsequent combinatorial theorems instantiate that lemma rather than
repeating sum rearrangements.

\subsection{Formalizing reflection}

Use a bit word $w<2^{2n}$.  Define
\begin{align*}
 \operatorname{up}(w,j)&=\sum_{i<j}\operatorname{bit}(w,i),\\
 \operatorname{down}(w,j)&=j\monus\operatorname{up}(w,j),\\
 \operatorname{dyck}(n,w)&\Longleftrightarrow
   \operatorname{up}(w,2n)=n\ \wedge\
   \sum_{j\le2n}[\operatorname{down}(w,j)>
                         \operatorname{up}(w,j)]=0.
\end{align*}
The first bad position is
\[
 \operatorname{firstBad}(w)=
 \min\{j\le2n:\operatorname{down}(w,j)>
                    \operatorname{up}(w,j)\},
\]
implemented by bounded fold with a found flag, not by unbounded minimization.
\texttt{reflectPrefix} toggles bits below that position.  Prove:
\begin{enumerate}[label=(\roman*),leftmargin=2em]
\item a bad balanced word maps to a word with $n+1$ up bits;
\item the first positive position of the image is the same cut;
\item applying the inverse prefix reflection recovers the input;
\item the maps preserve the bit-length bound.
\end{enumerate}
The generic bijection theorem then gives Equation~\eqref{eq:cat-difference}.
All four goals are bounded bit-count arithmetic and fit quantifier-free
induction on the prefix length.

\subsection{Hankel formalization without analytic dependencies}

Define \texttt{evenBallot r k} and \texttt{oddBallot r k} by the two binomial
differences.  The path-cut bijection first proves Equations
\eqref{eq:even-gram} and \eqref{eq:odd-gram}.  Instantiate ProveIt's
\texttt{det\_mul}, \texttt{det\_transpose}, and triangular determinant lemmas
for each encoded dimension $n$.  If the existing determinant interface is
awkward over natural numbers, work over the repository's integer-difference
ring and coerce all nonnegative entries.

For the third determinant, the only extra algebraic theorem is a rectangular
Cauchy--Binet identity.  A smaller alternative, useful for a first milestone,
is fraction-free elimination on
$[\Cat_{i+j+2}]$: the ballot matrix gives an explicit sequence of unimodular
row and column operations ending in
$\operatorname{diag}(1,\ldots,1,n+1)$.  Both routes are primitive recursive;
Cauchy--Binet better mirrors the human proof and generalizes to shifted
Catalan Hankel polynomials already studied elsewhere in ProveIt.

\subsection{Number-theoretic formalization}

The exact factorial module already imports the prime infrastructure.  Add a
fixed-base digit fold returning the triple
\((\text{digit sum},\text{carry count},\text{next carry})\).  Prove the digit
identity for addition by induction on the number of processed digits.  The
valuation formula then consists of:
\begin{enumerate}[label=(\arabic*),leftmargin=2em]
\item Legendre's bounded sum over powers $p^j$;
\item the finite digit-sum evaluation of that sum;
\item exact valuation of Equation~\eqref{eq:factorial};
\item the carry/digit identity;
\item specialization to $p=2$ and $p=5$.
\end{enumerate}
No unique factorization theorem stronger than the existing least-prime-divisor
induction is required.

\subsection{Formal series and real-code layers}

For formal series, use finite prefixes:
\begin{lstlisting}[language={}]
def coeff (a : Nat -> Int) (n : Nat) : Int := a n
def cauchy (a b : Nat -> Int) (n : Nat) : Int :=
  sumRange (n+1) (fun k => a k * b (n-k))
\end{lstlisting}
Every theorem quantifies over one requested coefficient $n$; there is no type
of completed power series in the object language.  Lagrange inversion can be
proved as a family of polynomial coefficient identities using rooted trees,
which is smaller than a formal derivative library.

For analytic claims, define a rational interval code
\begin{lstlisting}[language={}]
structure RealCode where
  lo hi : Nat -> RatCode
  ordered : forall q, lo q <= hi q
  width : forall q, q * (hi q - lo q) <= 1
  nested : ...
\end{lstlisting}
The ellipsis is not Lean code in the supplied scaffold; the report's
mathematical contract is monotone nesting.  Use Machin intervals for $\pi$,
integer bisection for square root, alternating intervals for logarithm, and
Taylor intervals for exponential.  The target asymptotic theorem is not a
statement about Lean's ambient \texttt{Real}; it is a formula saying that the
computed rational intervals obey Equation~\eqref{eq:leading-modulus}.

\subsection{Milestones and expected proof obligations}

A practical completion order is:
\begin{enumerate}[leftmargin=2em]
\item \textbf{Core}: closed forms, ratio, values, and Segner.
\item \textbf{Dyck/reflection}: one generic finite-bijection theorem and the
      path encodings.
\item \textbf{Finite identities}: all four MathWorld sums, product,
      hypergeometric termination, and $k$-fold convolution.
\item \textbf{Digits}: valuation, parity, final-digit rule, and Catalan primes.
\item \textbf{Hankel}: ballot Gram factorizations and all three determinants.
\item \textbf{Generalized}: rational, Fuss, Jonah, and super-Catalan results.
\item \textbf{Analysis}: boundary sum first (it needs no real library), then
      Bessel coefficient identities, integrals, Wallis, and Stirling.
\end{enumerate}

The finite milestones require only modules already present in the inspected
ProveIt dependency cone plus the new object encodings.  The analytic milestone
is larger, but every theorem statement and modulus has been reduced to
primitive-recursive rational data in this report.

\subsection{Status of the supplied Lean artifacts}

\texttt{CatalanExecutable.lean} contains executable natural-number definitions
and decidable small checks.  \texttt{CatalanProveItBlueprint.lean} records the
exact import, the checked theorem name, and definitions of the first target
terms.  The environment used to build this delivery did not contain a Lean
binary, so these new files are not claimed to have been compiled.  They use no
\texttt{sorry}, no new axioms, and no theorem declaration whose body is
missing.  The already existing ProveIt theorem \texttt{catalan} is separately
kernel checked in the repository.

\section{Coverage ledger, exclusions, and reproducibility}
\label{sec:coverage}

The detailed machine-readable ledger is
\texttt{reference/claim\_map.csv}.  The following table summarizes the
mathematical coverage.  ``Finite'' means an \ipra-ready proof by bounded
arithmetic or primitive-recursive bijection.  ``Effective'' means rational
Cauchy data and a primitive-recursive modulus have been supplied.  ``Target''
means that the claim is recorded but a dependency-heavy source-specific
formalization has not been completed.  ``Source correction means that the
requested page statement is false or misindexed as written; an explicit
counterexample and the corrected theorem are supplied.

\small
\begin{longtable}{@{}p{0.35\textwidth}p{0.15\textwidth}p{0.42\textwidth}@{}}
\toprule
\textbf{Claim family} & \textbf{Status} & \textbf{Location / method}\\
\midrule
\endfirsthead
\toprule
\textbf{Claim family} & \textbf{Status} & \textbf{Location / method}\\
\midrule
\endhead
Central-binomial, factorial, and difference formulas &
Checked base + finite & Section~\ref{sec:arithmetic}; ProveIt theorem
\texttt{catalan} plus reflection.\\
Odd-binomial and $n^{-1}\binom{2n}{n-1}$ forms & Finite &
Adjacent-binomial identities, Equations~\eqref{eq:odd-binomial}--\eqref{eq:n-binomial}.\\
Double factorial and product $\prod(4j-2)$ & Finite &
Equation~\eqref{eq:product}.\\
Generalized half-binomial formula & Finite rational &
Denominator-cleared expansion, Equation~\eqref{eq:half-binomial}.\\
Half-integer gamma formula & Effective / finite core &
Gamma recurrence reduces it to factorials, Equation~\eqref{eq:gamma}.\\
Terminating ${}_2F_1$ formula & Finite & Narayana/Vandermonde proof,
Equations~\eqref{eq:hypergeom}--\eqref{eq:narayana-term}.\\
Ratio recurrence & Finite & Theorem~\ref{thm:ratio-recurrence}.\\
Segner recurrence & Finite bijection & First-return decomposition,
Theorem~\ref{thm:segner}.\\
All six formula proofs displayed on Wikipedia & Finite & Generating functions,
reflection, Chung--Feller cut-and-swap, marked triangulations, balanced-string
convolution cancellation, and the cycle lemma.\\
MathWorld sums (9)--(12) & Finite & Two-colored Motzkin, weighted symmetry,
formal binomial transform, and Vandermonde square sum.\\
Riordan inverse transform & Finite bijection & Remove singleton blocks,
Equation~\eqref{eq:riordan-transform}.\\
Dyck words, balanced parentheses, grid paths, mountains & Finite bijections &
Sections~\ref{sec:arithmetic} and \ref{sec:combinatorics}.\\
Complete parenthesizations, binary parse trees, full binary trees & Finite
bijections & Recursive tree code, Equation~(3.2).\\
Rooted plane trees and planted trivalent trees & Finite bijections & Contour
traversal and left-child/right-sibling.\\
Polygon triangulations / Euler--Segner problem & Finite bijection & Root-edge
triangle split, Equation~(3.3).\\
Noncrossing handshakes / pair partitions & Finite bijection & First chord and
stack matching.\\
Noncrossing set partitions; $\Cat_n\le B_n$ & Finite & Root block formal
series and inclusion, Equations~\eqref{eq:NCseries}--\eqref{eq:bell-bound}.\\
$2$-by-$n$ standard Young tableaux & Finite bijection & Row-membership Dyck
word; hook specialization.\\
Restricted height sequences & Finite bijection & Equation~\eqref{eq:restricted-seq}.\\
All six length-three pattern classes & Finite & $231$ recursion, symmetries,
and finite RSK / square-sum identity.\\
Stack-sortable permutations & Finite & Recursive characterization equals
$231$ avoidance.\\
Staircase rectangle tilings & Finite bijection & Recursive antidiagonal split.\\
Positive integral frieze admissible diagonals & Target & Reduces to a known
triangulation correspondence, but exact frieze conventions are not encoded.\\
Oakley--Wisner flexagon configurations & Target & Specialized ``pat'' object
format is not reconstructed here.\\
Ordinary generating function & Finite coefficients + effective convergence &
$C=1+xC^2$, Equation~\eqref{eq:Cgf}.\\
Exponential generating function with $I_0-I_1$ & Effective & ODE and unique
coefficients, Equation~\eqref{eq:egf}.\\
Even Bessel generating function & Effective & Direct factorial coefficient
identity, Equation~\eqref{eq:even-egf}.\\
$\sum\Cat_n/4^n=2$ and sum from $n=1$ equals $1$ & Finite tail + effective &
Exact telescope, Theorem~\ref{thm:boundary-tail}.\\
Random-walk first-trap probabilities and recurrence & Finite + effective &
Equation~(4.7), tail at most $1/\sqrt{N+2}$.\\
Two integral representations & Effective & Beta reduction and rational
Riemann sums, Equation~\eqref{eq:integrals}.\\
Compact moment / normalized Hausdorff representation & Effective & Density
Equation~\eqref{eq:moment-density} and the rescaling $x=4y$.\\
Prime valuation formula & Finite & Legendre plus carry count,
Equation~\eqref{eq:valuation}.\\
Odd-prime digit rule & Finite automaton & Equations~(5.6)--(5.8).\\
Parity iff $n=2^k-1$ & Finite & Equation~\eqref{eq:v2} and
Equation~\eqref{eq:odd}.\\
Base-five last-digit criterion & Finite & Equations~\eqref{eq:base5-N}--\eqref{eq:base5-n}.\\
Only prime Catalans are $2,5$ & Finite & Prime-divisor bound and elementary
growth, Equation~(5.11).\\
Reported long finite exception searches & Reproducible only in configured
range & Criterion proved universally; source search ranges are not rebranded
as symbolic proofs.\\
Unshifted and once-shifted Hankel determinant $1$ & Finite & Ballot Gram
factorizations, Equation~\eqref{eq:hankel01}.\\
Twice-shifted Hankel determinant $n+1$ & Finite & Rectangular Gram,
Cauchy--Binet, unit maximal minors, Equation~\eqref{eq:hankel2}.\\
Two determinant sequences characterize Catalan & Finite & Southeast cofactor
induction, Equation~\eqref{eq:characterizing-Hankel}.\\
Ballot interpretation & Finite & Cycle lemma at excess one,
Equation~(7.1).\\
Catalan $k$-fold convolution & Finite & Ordered forests and Raney lemma,
Equation~\eqref{eq:kfold}.\\
Rational Catalan numbers & Finite & Ranked cyclic rotations,
Equation~\eqref{eq:rational-cat}.\\
Fuss--Catalan formula and $p$-ary trees & Finite & Polish words and cycle
lemma, Equation~\eqref{eq:fuss}.\\
$p$-ary operator associations, polygon dissections, $p$-good paths & Finite
bijections & Recursive $p$-ary tree models.\\
Generalized Jonah identity and explicit $d_{p,q}$ & Finite coefficients &
Formal reversion and Lagrange coefficient identity,
Equations~\eqref{eq:Dexplicit}--\eqref{eq:jonah}.\\
Generalized convolution (MathWorld (28)) & Finite coefficients & Product
identity, Equation~\eqref{eq:Dconvolution}.\\
Super-Catalan integrality and special cases & Finite & Prime-power floor
inequality, Equations~\eqref{eq:supercat}--(7.16).\\
General combinatorial interpretation of all super-Catalans & Open problem &
Explicitly not claimed.\\
Coxeter generalized Catalans, fully commutative elements, and diagram
deletion & Source corrections + type $A$ proof & Standard
$\operatorname{Cat}(A_n)=\Cat_{n+1}$; $B_2$ disproves the claimed general
fully-commutative identification; $A_2$ disproves the unweighted deletion
formula.  The corrected weighted type-$A$ recurrence is
Equation~\eqref{eq:coxeter-A-recurrence}; the all-types root-system development
remains a target.\\
Leading asymptotic & Effective & Wallis bound and modulus
Equation~\eqref{eq:leading-modulus}.\\
First three printed asymptotic terms & Effective & Stirling with remainder,
Equation~\eqref{eq:three-term-bound}.\\
Decimal digit ratio & Effective & Explicit error
Equation~\eqref{eq:digit-ratio}.\\
Complex-plane plots / unrestricted analytic continuation & Not a discrete
claim & No attempt to formalize graphical statements; the integer gamma values
needed here are covered.\\
Historical priority statements & Out of mathematical proof scope & Sources
are cited but historical claims are not re-proved in arithmetic.\\
\bottomrule
\end{longtable}
\normalsize

\subsection{Reproducibility artifacts}

The package contains:
\begin{itemize}[leftmargin=2em]
\item \texttt{code/verify\_catalan.py}, an exact-integer and rational checker
      for the closed forms, all finite sums, recurrences, valuations, Hankel
      determinants, generalized identities, and selected numerical data;
\item \texttt{verification/certificate.json}, generated by that script;
\item \texttt{lean/CatalanExecutable.lean}, an executable specification with
      small decidable checks;
\item \texttt{lean/CatalanProveItBlueprint.lean}, the object-language import
      and target-term scaffold;
\item \texttt{reference/source\_audit.md} and
      \texttt{reference/claim\_map.csv};
\item \texttt{MANIFEST.sha256}, checksums of delivered files.
\end{itemize}
The Python program is not an \ipra{} proof checker.  Its purpose is to catch
indexing and transcription errors and to generate finite certificates.  The
mathematical report supplies the uniform proofs.

\subsection{Conclusion}

The essential reason Catalan mathematics fits \ipra{} so well is that its
apparently analytic formulas rest on a small finite kernel: bounded words,
reflection, first return, Pascal arithmetic, and cyclic rotations.  Even the
boundary sum has an exact finite tail, the valuations are digit automata, and
the Hankel determinants are unit-triangular ballot factorizations.  The
remaining analytic notation can be compiled to rational intervals with
primitive-recursive moduli.  Thus a substantial portion of both requested
pages can be interpreted literally in intuitionistic primitive recursive
arithmetic rather than merely proved in a stronger metatheory and then called
``constructive.''

\begin{thebibliography}{99}

\bibitem{WikipediaCatalan}
Wikipedia contributors,
\emph{Catalan number}, revision accessed 16 November 2026.
\url{https://en.wikipedia.org/wiki/Catalan_number}.

\bibitem{MathWorldCatalan}
E. W. Weisstein,
\emph{Catalan Number}, MathWorld--A Wolfram Web Resource,
accessed 16 November 2026.
\url{https://mathworld.wolfram.com/CatalanNumber.html}.

\bibitem{ProveIt}
V. Reshetnikov,
\emph{ProveIt}, Git repository, inspected snapshot
\texttt{00bbe871268c453903499fb6afe07519e87d9631}.
\url{https://github.com/VladimirReshetnikov/ProveIt}.

\bibitem{StanleyEC2}
R. P. Stanley,
\emph{Enumerative Combinatorics, Volume 2},
2nd ed., Cambridge University Press, 2011.

\bibitem{ConcreteMath}
R. L. Graham, D. E. Knuth, and O. Patashnik,
\emph{Concrete Mathematics}, 2nd ed., Addison--Wesley, 1994.

\bibitem{KoshySalmassi}
T. Koshy and M. Salmassi,
Parity and primality of Catalan numbers,
\emph{College Mathematics Journal} 37 (2006), 52--53.

\bibitem{AlterKubota}
R. Alter and K. K. Kubota,
Prime and prime power divisibility of Catalan numbers,
\emph{Journal of Combinatorial Theory, Series A} 15 (1973), 243--256.

\bibitem{HiltonPedersen}
P. Hilton and J. Pedersen,
Catalan numbers, their generalization, and their uses,
\emph{Mathematical Intelligencer} 13 (1991), 64--75.

\bibitem{DvoretzkyMotzkin}
A. Dvoretzky and T. Motzkin,
A problem of arrangements,
\emph{Duke Mathematical Journal} 14 (1947), 305--313.

\bibitem{Robbins}
H. Robbins,
A remark on Stirling's formula,
\emph{American Mathematical Monthly} 62 (1955), 26--29.

\bibitem{SimionSchmidt}
R. Simion and F. W. Schmidt,
Restricted permutations,
\emph{European Journal of Combinatorics} 6 (1985), 383--406.

\bibitem{FominReading}
S. Fomin and N. Reading,
Generalized cluster complexes and Coxeter combinatorics,
\emph{International Mathematics Research Notices} 2005, no.~44,
2709--2757.

\bibitem{StembridgeFC}
J. R. Stembridge,
On the fully commutative elements of Coxeter groups,
\emph{Journal of Algebraic Combinatorics} 5 (1996), 353--385.

\end{thebibliography}


\end{document}
